% STAGE13-CHAPTER: V5-C06
% CHAPTER-SCHEMA-COVERAGE: CH-01,CH-02,CH-03,CH-04,CH-05,CH-06,CH-07,CH-08,CH-09,CH-10,CH-11,CH-12,CH-13,CH-14,CH-15,CH-16,CH-17,CH-18,CH-19,CH-20,CH-21,CH-22,CH-23,CH-24,CH-25,CH-26,CH-27,CH-28,CH-29,CH-30,CH-31,CH-32,CH-33,CH-34,CH-35,CH-36,CH-37
% CHAPTER-SCHEMA-STATUS: CH-01=passed; CH-02=passed; CH-03=passed; CH-04=passed; CH-05=passed; CH-06=passed; CH-07=passed; CH-08=passed; CH-09=passed; CH-10=passed; CH-11=passed; CH-12=passed; CH-13=passed; CH-14=passed; CH-15=passed; CH-16=passed; CH-17=passed; CH-18=passed; CH-19=passed; CH-20=passed; CH-21=passed; CH-22=passed; CH-23=passed; CH-24=passed; CH-25=passed; CH-26=passed; CH-27=passed; CH-28=passed; CH-29=passed; CH-30=passed; CH-31=passed; CH-32=passed; CH-33=passed; CH-34=passed; CH-35=passed; CH-36=passed; CH-37=passed
% CHAPTER-SCHEMA-EVIDENCE: CH-01=\label{struct:V5-C06-CH01}; CH-02=\label{struct:V5-C06-CH02}; CH-03=\label{struct:V5-C06-CH03}; CH-04=\label{struct:V5-C06-CH04}; CH-05=\label{struct:V5-C06-CH05}; CH-06=\label{struct:V5-C06-CH06}; CH-07=\label{struct:V5-C06-CH07}; CH-08=\label{struct:V5-C06-CH08}; CH-09=\label{struct:V5-C06-CH09}; CH-10=\label{struct:V5-C06-CH10}; CH-11=\label{struct:V5-C06-CH11}; CH-12=\label{struct:V5-C06-CH12}; CH-13=\label{struct:V5-C06-CH13}; CH-14=\label{struct:V5-C06-CH14}; CH-15=\label{struct:V5-C06-CH15}; CH-16=\label{struct:V5-C06-CH16}; CH-17=\label{struct:V5-C06-CH17}; CH-18=\label{struct:V5-C06-CH18}; CH-19=\label{struct:V5-C06-CH19}; CH-20=\label{struct:V5-C06-CH20}; CH-21=\label{struct:V5-C06-CH21}; CH-22=\label{struct:V5-C06-CH22}; CH-23=\label{struct:V5-C06-CH23}; CH-24=\label{struct:V5-C06-CH24}; CH-25=\label{struct:V5-C06-CH25}; CH-26=\label{struct:V5-C06-CH26}; CH-27=\label{struct:V5-C06-CH27}; CH-28=\label{struct:V5-C06-CH28}; CH-29=\label{struct:V5-C06-CH29}; CH-30=\label{struct:V5-C06-CH30}; CH-31=\label{struct:V5-C06-CH31}; CH-32=\label{struct:V5-C06-CH32}; CH-33=\label{struct:V5-C06-CH33}; CH-34=\label{struct:V5-C06-CH34}; CH-35=\label{struct:V5-C06-CH35}; CH-36=\label{struct:V5-C06-CH36}; CH-37=\label{struct:V5-C06-CH37}
% CHAPTER-MODULE-SEQUENCE: learning_navigation; strict_modeling; mathematical_development; multi_view_understanding; algorithm_engineering; examples_exercises; closure_self_check
% PROJECT_SPEC-V1.1-EXERCISE-LAYOUT: grouped; kept=16; source_group=4; supplement=12
% DERIVATION-CHECK: step-1; step-2; step-3
% FORMAL-RESULT-CHECK: results=7; proofs=7
% ALGORITHM-CHECK: algorithms=5; deterministic-contracts=5
% FIGURE-KNOWLEDGE-MAP: fig:V5-C06-dependency-graph -> SLM-20-00-001,PRE-TM-003,SLM-20-02-002,SLM-20-02-003
% FIGURE-KNOWLEDGE-MAP: fig:V5-C06-generative-process -> PRE-TM-003,SLM-20-02-004,SLM-20-01-005
% FIGURE-KNOWLEDGE-MAP: fig:V5-C06-plate-graph -> SLM-20-02-005
% FIGURE-KNOWLEDGE-MAP: fig:V5-C06-collapsed-gibbs-counts -> SLM-20-03-001,SLM-20-03-002,SLM-20-03-003,SLM-20-03-004,SLM-20-03-005,SLM-20-02-008
% FIGURE-KNOWLEDGE-MAP: fig:V5-C06-elbo-geometry -> SLM-20-04-001,SLM-20-04-002,SLM-20-04-003,SLM-20-03-006
% FIGURE-KNOWLEDGE-MAP: fig:V5-C06-mean-field-graph -> SLM-20-04-004
% FIGURE-KNOWLEDGE-MAP: fig:V5-C06-variational-updates -> SLM-20-04-004,SLM-20-04-005,SLM-20-04-006,SLM-20-05-001
% FIGURE-KNOWLEDGE-MAP: fig:V5-C06-method-comparison -> SLM-20-03-001,SLM-20-04-001,SLM-20-04-006,SLM-20-05-001
% REPRODUCTION-NODE: V5-C06-S02 -> PRE-TM-003
% SOURCE-BOUNDARY-DIFFERENCE: generic Dirichlet--multinomial conjugacy belongs to V5-C05; this chapter cites and specializes it without re-proving the generic result.

\chapter{潜在狄利克雷分配}\label{chap:V5-C06}
\index{潜在狄利克雷分配}
\index{LDA|see{潜在狄利克雷分配}}
\index{主题模型}

% KNOWLEDGE-PLAN: KN-V5-C35-PREREQUISITE-001 L1
\paragraph{读前自检：本章地图：潜在狄利克雷分配。}潜在狄利克雷分配章地图把“词袋语料$\to$LDA生成模型$\to$折叠Gibbs/变分推断$\to$留出评价”作为主线；请为其中首条箭头补出必要条件，并检查终点仍落在本章声明的输出域。\par

\SLChapterOpening{词袋语料$\to$LDA生成模型$\to$折叠Gibbs/变分推断$\to$留出评价}
{怎样从概率图完整推导LDA联合分布，并合法消元或构造变分近似？}
{读完可推出折叠计数条件、KL--ELBO恒等式和可归一化的坐标更新。}
{条件独立、Dirichlet共轭、Gibbs抽样、KL散度、Jensen不等式、EM与Newton法。}
{先掌握\GlobalChapterRef{34}{5}{5}的多元Beta积分和\GlobalChapterRef{33}{5}{4}的坐标Gibbs，再进入生成过程。}

\phantomsection\label{struct:V5-C06-CH01}
% KNOWLEDGE-PLAN: KN-V5-C35-PREREQUISITE-002 L1
\paragraph{读前自检：本章可检验学习目标。}潜在狄利克雷分配学习目标固定核验“把语料、文档、主题和单词写成维数明确的随机变量，并从概率图读出完整联合分布”；请实际完成“把语料、文档、主题和单词写成维数明确的随机变量，并从概率图读出完整联合分布”，并用“中间结果逐项包含首目标“把语料、文档、主题和单词写成维数明确的随机变量，并从概率图读出完整联合分布”声明的对象、条件与结论”判定通过或失败。\par

\chaptergoals{
\begin{itemize}[leftmargin=2em]
  \item 把语料、文档、主题和单词写成维数明确的随机变量，并从概率图读出完整联合分布；
  \item 解释文档--主题分布与主题--词分布的两层混合，以及词袋交换性所允许和不允许的结论；
  \item 由\DirichletMultinomial{}积分逐步推出折叠联合分布、留一计数和Gibbs满条件分布；
  \item 从KL--ELBO恒等式推出平均场坐标更新、主题--词参数更新和超参数Newton更新；
  \item 按严格输入输出契约复现文本生成、折叠Gibbs、通用变分EM、局部变分推断和完整变分EM；
  \item 用留出似然、困惑度、稳定性和资源复杂度评价主题模型，而不把训练目标误当成泛化能力。
\end{itemize}
}

\phantomsection\label{struct:V5-C06-CH02}
% KNOWLEDGE-PLAN: KN-V5-C35-PREREQUISITE-003 L1
\paragraph{读前自检：最短先修。}LDA最短先修要求能归一化Dirichlet--类别模型并区分KL、Jensen与Gibbs条件更新；请对$K$维正参数$\alpha$写出$E\theta_k=\alpha_k/\alpha_0$，核对各主题权重和为1。\par

\begin{prerequisitebox}
本章直接使用离散随机变量与条件独立、Bayes公式、KL散度、Jensen不等式、
\DirichletMultinomial{}共轭、Gibbs抽样、EM算法和Newton法。尤其需要熟悉\GlobalChapterRef{34}{5}{5}的结论：
Dirichlet先验与类别计数相乘后，后验参数按分量加上计数；这一通用共轭结论在本章只作局部调用，不重复证明。读者还应能区分概率向量、计数向量和独热向量。
\end{prerequisitebox}

\phantomsection\label{struct:V5-C06-CH03}
% KNOWLEDGE-PLAN: KN-V5-C35-PREREQUISITE-004 L1
\paragraph{读前自检：先修自检。}潜在狄利克雷分配先修自检固定核验“能否把$p(a,b,c)=p(a)p(b\mid a)p(c\mid b)$画成有向图，并从图中说明$c\perp a\mid b$”；请把$p(a,b,c)=p(a)p(b\mid a)p(c\mid b)$画成有向图，并从图中说明$c\perp a\mid b$并写出中间结果，再按“$p(a,b,c)=p(a)p(b\mid a)p(c\mid b)$的计算或证明结果满足首项所列等式、定义域或判断”明确判为通过或失败。\par

\begin{precheckbox}
\begin{enumerate}[leftmargin=2.2em]
  \item 能否把$p(a,b,c)=p(a)p(b\mid a)p(c\mid b)$画成有向图，并从图中说明$c\perp a\mid b$？
  \item 若$\boldsymbol\theta\sim\operatorname{Dir}(\boldsymbol\alpha)$且类别计数为$\boldsymbol n$，能否写出积分
  $\int\prod_k\theta_k^{n_k}p(\boldsymbol\theta\mid\boldsymbol\alpha)\,\mathrm d\boldsymbol\theta$的Beta函数比值？
  \item 能否说明Gibbs抽样中“先删当前状态再计算条件概率”的原因？
  \item 能否由$\log p(x)=\mathcal L(q)+\KL(q\|p)$判断最大化ELBO与最小化KL的关系？
  \item 能否在单纯形约束下用拉格朗日乘子把一组正权重归一化？
\end{enumerate}
第1项不足时回看概率图与条件独立；第2项不足时回看\GlobalChapterRef{34}{5}{5}的\DirichletMultinomial{}共轭；
第3项不足时回看Gibbs抽样；第4--5项不足时回看变分推断、EM和约束优化。
\end{precheckbox}

\phantomsection\label{struct:V5-C06-CH04}
% KNOWLEDGE-PLAN: KN-V5-C35-PREREQUISITE-005 L1
\paragraph{读前自检：依赖路线。}潜在狄利克雷分配把“所引章节的\DirichletMultinomial{}共轭与所引章节的Gibbs坐标更新 $\to$ 本章两条推断路线”接到“有限类别与独热编码$\to$文档、主题和单词变量”；请用$\to$补全这一步的等式或判据，再检查两端维数、支持或对象类型。\par

\begin{dependencybox}
\GlobalChapterRef{34}{5}{5}的\DirichletMultinomial{}共轭与\GlobalChapterRef{33}{5}{4}的Gibbs坐标更新 $\to$ 本章两条推断路线；
有限类别与独热编码$\to$文档、主题和单词变量；
\DirichletMultinomial{}共轭$\to$积分消去$\boldsymbol\theta,\boldsymbol\varphi$
$\to$折叠Gibbs；
KL散度与Jensen不等式$\to$ELBO$\to$平均场坐标更新；
EM与Newton法$\to$全局主题参数和超参数学习；
留出预测$\to$困惑度与模型比较。图\ref{fig:V5-C06-dependency-graph}只画这些直接依赖，不把诊断性远缘关系冒充前置课程。
\end{dependencybox}
% v2.3.1 figure UID: FIG-P680-01
% Proposed post-v2.3.1 polish for FIG-P680-01: protect key-label size and stroke hierarchy.
\tikzset{slfig-FIG-P680-01/.style={font=\fontsize{9.4pt}{11.4pt}\selectfont,every path/.append style={line cap=round,line join=round}}}
% v2.1.0 priority-figure QA: 9.8pt labels, strengthened strokes, semantic direct labels.
\pgfkeysifdefined{/tikz/alt}{}{\tikzset{alt/.code={}}}
\begin{figure}[htbp]
\centering
\begin{tikzpicture}[slfig-FIG-P680-01,slfig high,>=Stealth,
  every node/.style={font=\fontsize{9.4pt}{11.4pt}\selectfont,align=center},
  core/.style={draw=SLBlue,fill=SLBlueBG,line width=.80pt,rounded corners=2.5pt,minimum width=54mm,minimum height=12mm,inner xsep=3mm,inner ysep=1.4mm},
  leftmodel/.style={draw=SLTeal,fill=SLTealBG,line width=.76pt,rounded corners=2.5pt,minimum width=46mm,minimum height=12mm,inner xsep=3mm,inner ysep=1.4mm},
  rightmodel/.style={draw=SLGold,fill=SLGoldBG,line width=.76pt,rounded corners=2.5pt,minimum width=46mm,minimum height=12mm,inner xsep=3mm,inner ysep=1.4mm},
  rowlabel/.style={font=\fontsize{9.8pt}{11.8pt}\selectfont\bfseries,text=SLGray,anchor=east},
  warning/.style={draw=SLRed,fill=SLRedBG,text=SLRed,line width=.76pt,rounded corners=2.5pt,minimum width=116mm,minimum height=9mm,inner xsep=3mm,inner ysep=1.5mm,font=\fontsize{9.2pt}{11.2pt}\selectfont},
  alt={对象—关系—结论：本章从共享的文档—主题—单词条件结构分出两个模型目标：完整Bayes LDA由折叠Gibbs推断，点参数LDA变体由平均场变分EM估计；箭头表示学习依赖，不表示两个后验相同}]
\node[core] (shared) at (0,2.25) {共享“文档--主题--单词”条件结构\\两层混合、条件独立与词袋计数};
\node[rowlabel] at (-5.85,.65) {模型目标};
\node[leftmodel] (full) at (-3.10,.65) {完整 Bayes LDA\\$\theta,\varphi$ 均为随机变量};
\node[rightmodel] (point) at (3.10,.65) {点参数 LDA 变体\\$\varphi$ 作为待估参数};
\node[rowlabel] at (-5.85,-.95) {推断方法};
\node[leftmodel,fill=white] (gibbs) at (-3.10,-.95) {折叠 Gibbs\\积分消去 $\theta,\varphi$};
\node[rightmodel,fill=white] (vem) at (3.10,-.95) {平均场变分 EM\\ELBO 坐标上升};
\draw[-{Stealth[length=2.1mm,width=1.55mm]},draw=SLTeal,line width=.95pt]
  ([xshift=-7mm]shared.south) to[bend right=12] (full.north);
\draw[-{Stealth[open,length=2.1mm,width=1.55mm]},draw=SLGold,line width=.95pt,densely dashed]
  ([xshift=7mm]shared.south) to[bend left=12] (point.north);
\draw[-{Stealth[length=1.9mm,width=1.4mm]},draw=SLTeal,line width=.80pt] (full.south)--(gibbs.north);
\draw[-{Stealth[open,length=1.9mm,width=1.4mm]},draw=SLGold,line width=.80pt,densely dashed] (point.south)--(vem.north);
\node[warning] at (0,-2.40) {模型与后验不同，结果不可只按算法名直接比较};
\end{tikzpicture}
\captionsetup{width=.94\linewidth}
\caption{本章从共享的文档--主题--单词条件结构分出两个模型目标：完整Bayes LDA由折叠Gibbs推断，点参数LDA变体由平均场变分EM估计；箭头表示学习依赖，不表示两个后验相同}
\label{fig:V5-C06-dependency-graph}
\end{figure}


\begin{keypointbox}[title=数据方向桥：LDA语料仍按列放文档]
若把词元压缩成计数矩阵，本章采用$X_{\rm wd}\in\mathbb N_0^{V\times M}$：词表项在行，文档在列。与全书默认行样本表示的关系是
\[
X_{{\rm doc},{\rm row}}=X_{\rm wd}^{\mathsf T}
\in\mathbb N_0^{M\times V}.
\]
例如$V=1000,M=200$时，两种维数为$1000\times200$与$200\times1000$。LDA内部的$\theta_m$仍对应一篇文档，$\varphi_k$仍对应一个主题；矩阵转置不交换这两个概率对象的角色。
\end{keypointbox}

\phantomsection\label{struct:V5-C06-CH05}
% STAGE19-SYMBOL-INDEX-BEGIN: V5-C06
\index[symbols]{m-v-k--s0a553ec1313c@$M,V,K$：文档数、词表大小、主题数；角色：三个基本规模参数}
\index[symbols]{n-m-i--sab7ca9215267@$N_m,I$：第$m$篇长度、总词元数；角色：计算复杂度的样本规模}
\index[symbols]{w-mn--se5a53a82594a@$w_{mn}$：第$m$篇第$n$个词元；角色：观测变量}
\index[symbols]{z-mn--s32a8da2181af@$z_{mn}$：词元的主题指派；角色：离散隐变量}
\index[symbols]{minus-boldsymbol-minus-minus-theta-minus-m--sf7d47f7f095a@$\boldsymbol\theta_m$：文档--主题概率向量；角色：文档级隐变量}
\index[symbols]{minus-boldsymbol-minus-minus-varphi-minus-k--sb2217db38d80@$\boldsymbol\varphi_k$：主题--词概率向量；角色：主题级参数或隐变量}
\index[symbols]{minus-boldsymbol-minus-minus-alpha-minus--s4f106bbffd01@$\boldsymbol\alpha$：文档主题先验；角色：Dirichlet超参数}
\index[symbols]{minus-boldsymbol-minus-minus-beta-minus--sb5087e9bdbf6@$\boldsymbol\beta$：主题词先验；角色：Dirichlet超参数}
\index[symbols]{n-mk-n-kv--s0a8553e38c22@$n_{mk},n_{kv}$：文档--主题、主题--词计数；角色：折叠推断充分统计量}
\index[symbols]{minus-boldsymbol-minus-minus-gamma-minus-m--sded0b7f677e5@$\boldsymbol\gamma_m$：$q(\boldsymbol\theta_m)$参数；角色：局部变分参数}
\index[symbols]{minus-boldsymbol-minus-minus-eta-minus-mn--s055f19bb8ebf@$\boldsymbol\eta_{mn}$：$q(z_{mn})$参数；角色：主题责任度}
\index[symbols]{minus-psi-minus-minus-psi-minus-u0027--s3f066ea944e0@$\psi,\psi'$：digamma、trigamma；角色：Dirichlet对数矩与Newton步}
% STAGE19-SYMBOL-INDEX-END: V5-C06
% KNOWLEDGE-PLAN: KN-V5-C35-PREREQUISITE-006 L1
\paragraph{读前自检：本章符号表。}潜在狄利克雷分配符号表规定$M,V,K$、$N_m,I$的类型或范围；请把$M,V,K$代入紧随其后的定义关系，逐项核对输入形状、输出形状与取值域。\par

\begin{symboltablebox}
\small
\begin{longtable}{@{}p{0.17\textwidth}p{0.25\textwidth}p{0.22\textwidth}p{0.27\textwidth}@{}}
\toprule
符号 & 对象与读法 & 维数/范围 & 角色 \\
\midrule
$M,V,K$ & 文档数、词表大小、主题数 & 正整数 & 三个基本规模参数 \\
$N_m,I$ & 第$m$篇长度、总词元数 & $I=\sum_mN_m$ & 计算复杂度的样本规模 \\
$w_{mn}$ & 第$m$篇第$n$个词元 & $\{1,\ldots,V\}$ & 观测变量 \\
$z_{mn}$ & 词元的主题指派 & $\{1,\ldots,K\}$ & 离散隐变量 \\
$\boldsymbol\theta_m$ & 文档--主题概率向量 & $\Delta^{K-1}$ & 文档级隐变量 \\
$\boldsymbol\varphi_k$ & 主题--词概率向量 & $\Delta^{V-1}$ & 主题级参数或隐变量 \\
$\boldsymbol\alpha$ & 文档主题先验 & $\mathbb R_{>0}^{K}$ & Dirichlet超参数 \\
$\boldsymbol\beta$ & 主题词先验 & $\mathbb R_{>0}^{V}$ & Dirichlet超参数 \\
$n_{mk},n_{kv}$ & 文档--主题、主题--词计数 & 非负整数 & 折叠推断充分统计量 \\
$\boldsymbol\gamma_m$ & $q(\boldsymbol\theta_m)$参数 & $\mathbb R_{>0}^{K}$ & 局部变分参数 \\
$\boldsymbol\eta_{mn}$ & $q(z_{mn})$参数 & $\Delta^{K-1}$ & 主题责任度 \\
$\psi,\psi'$ & digamma、trigamma & 标量函数 & Dirichlet对数矩与Newton步 \\
\bottomrule
\end{longtable}
为使低维数值例的式面更紧凑，下文在等号右侧直接列出全部坐标时采用如下显式排版约定：
$\boldsymbol\alpha,\boldsymbol\beta,\boldsymbol\theta,\boldsymbol\varphi,
\boldsymbol\eta,\boldsymbol\gamma$表示向量，而带下标的$\alpha_k,\beta_v$等表示分量；括号中的完整数列按符号表给出的维数解释为对应向量。
该约定不改变对象的维数、归一化约束或推导含义。
\end{symboltablebox}

% KNOWLEDGE-PLAN: KN-V5-C35-CONCEPT-002 L1
\paragraph{读前自检：潜在狄利克雷分配。}LDA把文档表示成主题混合$\theta_m$，再给每个词位置分配主题$z_{mn}$并生成$w_{mn}$；请沿“文档--主题混合--位置主题--单词”写出一项条件概率，核对“分配”指位置级主题指派。\par
% KNOWLEDGE-PLAN: KN-V5-C35-CONCEPT-003 L1
\paragraph{读前自检：LDA生成过程与概率图。}LDA生成过程先采样$\varphi_k\sim\operatorname{Dir}(\beta)$和$\theta_m\sim\operatorname{Dir}(\alpha)$，再采样$z_{mn},w_{mn}$；请按箭头顺序写出联合分布的四类因子并核对板内索引。\par
% KNOWLEDGE-PLAN: KN-V5-C35-CONCEPT-004 L1
\paragraph{读前自检：LDA文档—主题—词模型。}LDA文档--主题--词模型满足$P(w_{mn}=v\mid\theta_m,\varphi)=\sum_k\theta_{mk}\varphi_{kv}$；请对$v$求和，核对主题混合与主题词分布均归一时得到1。\par

\SLReviewSection{文档—主题—单词关系}{sec:V5-C06-S01}
\knowledgeanchor{SLM-20-00-001}{潜在狄利克雷分配}
\knowledgeanchor{PRE-TM-003}{LDA生成过程与概率图}
\knowledgeanchor{SLM-20-02-002}{LDA文档—主题—词模型}
\knowledgeanchor{SLM-20-02-003}{LDA的混合主题思想}

\phantomsection\label{struct:V5-C06-CH06}
\textbf{问题提出。}一个文档通常不只谈一个主题，一个主题也不会只使用一个词。
若直接给每篇文档估计一个$V$维词频向量，短文档会非常稀疏，文档之间也无法共享统计强度。
LDA在二者之间加入$K$个潜在主题：每篇文档先混合主题，每个词位置再选择主题，最后由所选主题产生单词。于是
\[
\text{文档}\longrightarrow\text{主题混合}\longrightarrow
\text{逐位置主题}\longrightarrow\text{单词}.
\]
两层Dirichlet先验分别平滑文档--主题分布和主题--词分布。学习任务不是为文档贴唯一标签，而是从只观察到的词序列反推主题词义、文档主题比例和逐词主题指派。

\begin{notationbox}
“潜在”指$z_{mn}$和$\boldsymbol\theta_m$未被直接观察；“狄利克雷”指概率向量的先验；
“分配”指每个词位置的主题指派。LDA不是线性判别分析，也不等于把文档硬聚成一个主题。
\end{notationbox}


\phantomsection\label{struct:V5-C06-CH07}
% KNOWLEDGE-PLAN: KN-V5-C35-DEFINITION-001 L1
\paragraph{读前自检：潜在狄利克雷分配模型。}LDA给定$M$篇文档、$K$个主题和$V$词表，要求$\theta_m\in\Delta^{K-1}$、$\varphi_k\in\Delta^{V-1}$；请核对两向量维数与归一化，并标出$z_{mn}$和$w_{mn}$的取值集合。\par

\begin{definition}[潜在狄利克雷分配模型]\label{def:V5-C06-lda}
给定词表$\mathcal W=\{1,\ldots,V\}$、主题数$K$和文档数$M$。
对每个主题$k$，取
$\boldsymbol\varphi_k\sim\operatorname{Dir}(\boldsymbol\beta)$；
对每篇文档$m$，取
$\boldsymbol\theta_m\sim\operatorname{Dir}(\boldsymbol\alpha)$。
对文档$m$的每个位置$n$，依次生成
\[
z_{mn}\mid\boldsymbol\theta_m\sim\operatorname{Cat}(\boldsymbol\theta_m),
\qquad
w_{mn}\mid z_{mn}=k,\boldsymbol\varphi
\sim\operatorname{Cat}(\boldsymbol\varphi_k).
\]
所有$\alpha_k>0$、$\beta_v>0$。本章明确区分两个相关但目标不同的模型：
\begin{enumerate}[leftmargin=2em]
  \item \textbf{完整Bayes LDA：}$\boldsymbol\varphi_k\sim\operatorname{Dir}(\boldsymbol\beta)$，折叠Gibbs积分消去$\boldsymbol\theta,\boldsymbol\varphi$；
  \item \textbf{点参数LDA变体：}$\boldsymbol\varphi_k\in\Delta^{V-1}$是待估参数，不对$\boldsymbol\varphi_k$放置Dirichlet先验，变分EM最大化该参数模型的ELBO。
\end{enumerate}
两者共享条件似然$p(w_{mn}\mid z_{mn},\boldsymbol\varphi)$和文档先验
$\boldsymbol\theta_m\sim\operatorname{Dir}(\boldsymbol\alpha)$，但不是同一个后验，也不能把
$\boldsymbol\beta$平滑项临时加入点参数变体后仍声称优化同一目标。
\end{definition}

\SLReadTranslation{首次阅读固定走“完整Bayes模型$\to$积分消去$\boldsymbol\theta,\boldsymbol\varphi$$\to$折叠Gibbs满条件$\to$减—算—抽—加”这一条完整路线。第\ref{sec:V5-C06-S04}节的变分EM是对点参数变体的进阶对照：它优化ELBO，不是对前一条完整Bayes后验换一种算法。}

% KNOWLEDGE-PLAN: KN-V5-C35-DEFINITION-002 L1
\paragraph{读前自检：模型定义。}LDA模型定义区分对$\theta,\varphi$积分的完整Bayes模型与点参数变体；请写出后验$p(z,\theta,\varphi\mid w,\alpha,\beta)$的分子，核对折叠Gibbs与点参数变分EM针对的未知对象不同。\par
% KNOWLEDGE-PLAN: KN-V5-C35-ALGORITHM_IDEA-001 L1
\paragraph{读前自检：LDA的文本生成算法。}LDA文本生成算法以$M,K,V,N_m,\alpha,\beta$为输入；请先采样一列$\varphi_k$和一篇文档的$\theta_m$，再采样一个$z_{mn},w_{mn}$，核对概率向量归一并在全部计划对象生成后停止。\par
% KNOWLEDGE-PLAN: KN-V5-C35-CONCEPT-006 L1
\paragraph{读前自检：概率图模型。}LDA概率图在给定$\theta_m$与全部$\varphi_k$时，使同一文档各位置条件独立；请写出$P(z_{mn},w_{mn}\mid\theta_m,\varphi)$，核对跨位置共享$\theta_m$仍会产生边缘相关。\par
% KNOWLEDGE-PLAN: KN-V5-C35-PREREQUISITE-007 L1
\paragraph{读前自检：随机变量序列的可交换性。}LDA文档内词位置在积分掉共享$\theta_m$后可交换但一般不独立；请置换两个位置并比较条件联合乘积，核对乘积因子次序改变而数值不变。\par
% KNOWLEDGE-PLAN: KN-V5-C35-FORMULA-001 L1
\paragraph{读前自检：概率公式。}LDA完整联合分布由主题词先验、文档主题先验、主题指派和发词因子相乘；请对一个位置写出$\theta_{m,z_{mn}}\varphi_{z_{mn},w_{mn}}$，核对所有索引落在$M,K,V$声明范围。\par

\SLTeachSection{LDA生成过程与概率图}{sec:V5-C06-S02}
\knowledgeanchor{SLM-20-02-004}{模型定义}
\knowledgeanchor{SLM-20-01-005}{LDA的文本生成算法}
\knowledgeanchor{SLM-20-02-005}{概率图模型}
\knowledgeanchor{SLM-20-02-006}{随机变量序列的可交换性}
\knowledgeanchor{SLM-20-02-007}{概率公式}

\phantomsection\label{struct:V5-C06-CH08}
对完整Bayes LDA，联合分布分解为
\begin{align}
&p(\boldsymbol w,\boldsymbol z,\boldsymbol\theta,
\boldsymbol\varphi\mid\boldsymbol\alpha,\boldsymbol\beta)
\label{eq:V5-C06-joint}\\
&\quad=
\prod_{k=1}^{K}p(\boldsymbol\varphi_k\mid\boldsymbol\beta)
\prod_{m=1}^{M}\left[
p(\boldsymbol\theta_m\mid\boldsymbol\alpha)
\prod_{n=1}^{N_m}
p(z_{mn}\mid\boldsymbol\theta_m)
p(w_{mn}\mid z_{mn},\boldsymbol\varphi)
\right].\notag
\end{align}
点参数变体则以$\boldsymbol\varphi$为给定参数，其联合分布为
\begin{equation}
p_{\mathrm{par}}(\boldsymbol w,\boldsymbol z,\boldsymbol\theta
\mid\boldsymbol\alpha,\boldsymbol\varphi)
=\prod_{m=1}^{M}\left[
p(\boldsymbol\theta_m\mid\boldsymbol\alpha)
\prod_{n=1}^{N_m}p(z_{mn}\mid\boldsymbol\theta_m)
p(w_{mn}\mid z_{mn},\boldsymbol\varphi)
\right].
\label{eq:V5-C06-parameterized-joint}
\end{equation}
板块图中，$N_m$板块嵌在$M$板块内，表示每篇文档的每个词位置都重复生成；
$K$板块表示主题词向量重复$K$次。图\ref{fig:V5-C06-plate-graph}用统一符号展示这一生成依赖。
% v2.6.0 figure UID: FIG-P683-01
% Semantic lock: 18 node / 5 draw, N_m nested in M, K separate; no resizebox.
\tikzset{slfig-FIG-P683-01/.style={font=\fontsize{9.4pt}{11.3pt}\selectfont,every path/.append style={line cap=round,line join=round}}}
\pgfkeysifdefined{/tikz/alt}{}{\tikzset{alt/.code={}}}
\begin{figure}[htbp]
\centering
\begin{tikzpicture}[slfig-FIG-P683-01,slfig high,>=Stealth,
  every node/.style={font=\fontsize{9.4pt}{11.3pt}\selectfont},
  latent/.style={circle,draw=SLTeal,fill=SLTealBG,line width=.80pt,minimum size=8.2mm,inner sep=0pt,font=\fontsize{9.6pt}{11.5pt}\selectfont},
  hyper/.style={draw=none,fill=none,inner sep=1pt,text=SLGray,font=\fontsize{9.6pt}{11.5pt}\selectfont},
  obs/.style={circle,draw=SLBlue,fill=SLBlue,line width=.85pt,text=white,minimum size=8.2mm,inner sep=0pt,font=\bfseries\fontsize{9.6pt}{11.5pt}\selectfont},
  edge/.style={-{Stealth[length=2.1mm,width=1.5mm]},draw=SLBlue,line width=.95pt},
  plate/.style={draw=SLGray!62,line width=.66pt,rounded corners=2.2pt,dash pattern=on 2.6pt off 1.8pt},
  plabel/.style={text=SLGray,font=\fontsize{9.2pt}{11.0pt}\selectfont,inner sep=0pt},
  legendtext/.style={anchor=west,font=\fontsize{9.2pt}{11.0pt}\selectfont,inner sep=0pt},
  alt={对象—关系—结论：完整Bayes LDA盘式图把超参数、潜变量和观测变量分开：每篇文档共享一个主题比例，每个词位拥有一个主题指派，所有文档共享带Dirichlet先验的主题词分布；盘框标明重复次数，箭头只表示条件依赖方向}]
\node[hyper] (alpha) at (-5.95,1.25) {$\alpha$};
\node[latent] (theta) at (-3.05,1.25) {$\theta_m$};
\node[latent] (z) at (-.65,1.25) {$z_{mn}$};
\node[obs] (w) at (1.65,1.25) {$w_{mn}$};
\node[hyper] (beta) at (-5.95,-2.10) {$\beta$};
\node[latent] (phi) at (-.65,-2.10) {$\varphi_k$};

\node[plate,fit=(z)(w),inner xsep=4.6mm,inner ysep=4.4mm] (nplate) {};
\node[plate,fit=(theta)(nplate),inner xsep=5.2mm,inner ysep=6.8mm] (mplate) {};
\node[plate,fit=(phi),inner xsep=4.5mm,inner ysep=3.5mm] (kplate) {};

\draw[edge] (alpha) -- (theta);
\draw[edge] (theta) -- (z);
\draw[edge] (z) -- (w);
\draw[edge] (beta) -- (phi);
\draw[edge] (phi) -- (w);

\node[plabel,anchor=south east] at ($(nplate.north east)+(-1.5mm,1.55mm)$) {$N_m$ 个词位};
\node[plabel,anchor=north east] at ($(mplate.south east)+(0,-1.9mm)$) {$M$ 篇文档};
\node[plabel,anchor=north east] at ($(kplate.south east)+(0,-1.9mm)$) {$K$ 个主题};
\node[obs,minimum size=6.4mm] at (5.45,.95) {};
\node[legendtext] at (6.25,.95) {观测变量};
\node[latent,minimum size=6.4mm] at (5.45,0) {};
\node[legendtext] at (6.25,0) {潜变量};
\node[hyper] at (5.45,-.95) {$\alpha,\beta$};
\node[legendtext] at (6.25,-.95) {超参数（plate 外）};
\end{tikzpicture}
\captionsetup{width=.94\linewidth}
\caption{完整Bayes LDA盘式图把超参数、潜变量和观测变量分开：每篇文档共享一个主题比例，每个词位拥有一个主题指派，所有文档共享带Dirichlet先验的主题词分布；盘框标明重复次数，箭头只表示条件依赖方向}
\label{fig:V5-C06-plate-graph}
\end{figure}


\phantomsection\label{struct:V5-C06-CH09}
\textbf{学习策略。}已知$\boldsymbol w$后，目标是后验
\[
p(\boldsymbol z,\boldsymbol\theta,\boldsymbol\varphi
\mid\boldsymbol w,\boldsymbol\alpha,\boldsymbol\beta).
\]
对$\boldsymbol z$求和并对两组概率向量积分会产生指数规模耦合，通常不能精确计算。
折叠Gibbs在完整Bayes模型中积分消去共轭概率向量，再对$\boldsymbol z$抽样；
变分EM在式\eqref{eq:V5-C06-parameterized-joint}的点参数变体中，以可计算的平均场分布近似
$p_{\mathrm{par}}(\boldsymbol\theta,\boldsymbol z\mid\boldsymbol w,
\boldsymbol\alpha,\boldsymbol\varphi)$并交替更新局部变分参数与全局点参数。
前者提供完整Bayes后验的相关样本，后者提供另一个明确参数模型的确定性局部近似。

\begin{proposition}[条件独立生成蕴含文档内交换性]\label{prop:V5-C06-exchangeable}
在给定$\boldsymbol\theta_m$和$\boldsymbol\varphi$时，文档$m$各词位置独立同分布，且
\[
p(w_{mn}=v\mid\boldsymbol\theta_m,\boldsymbol\varphi)
=\sum_{k=1}^{K}\theta_{mk}\varphi_{kv}.
\]
因此任意置换词位置不改变条件联合概率；积分消去
$\boldsymbol\theta_m$后，词位置仍交换，但一般不再边缘独立。
\end{proposition}
% KNOWLEDGE-PLAN: KN-V5-C35-PROOF-001 L1
\paragraph{读前自检：证明：条件独立生成蕴含文档内交换性。}LDA文档内词元在给定$\boldsymbol\theta_m$时的联合概率为$\prod_n\sum_k\theta_{mk}\varphi_{k,w_{mn}}$；请置换位置指标核对乘积不变，再指出边缘化后各位置仍因共享$\boldsymbol\theta_m$而不独立。\par

\begin{proof}
对任意位置置换$\pi$，条件联合概率是
\[
\prod_{n=1}^{N_m}\sum_k\theta_{mk}\varphi_{k,w_{mn}},
\]
乘法交换律使其等于置换后的乘积。再对
$p(\boldsymbol\theta_m\mid\boldsymbol\alpha)$积分不会破坏置换不变性。
但两个位置共享同一个随机$\boldsymbol\theta_m$；边缘化后它们通过该共同原因相关，所以不能由交换性推出独立性。
\end{proof}

\phantomsection\label{struct:V5-C06-CH10}
\textbf{生成算法。}图\ref{fig:V5-C06-generative-process}以三条泳道展示完整Bayes LDA 的四步祖先采样；$n=1,\ldots,N_m$嵌套在每个$m$内，侧注只用于区分点参数变体。
% v2.7.0 figure UID: FIG-P684-01
% Three-lane ancestral sampler; intentionally not a duplicate full plate graph.
\tikzset{slfig-FIG-P684-01/.style={
  font=\fontsize{9.6pt}{11.5pt}\selectfont,
  every path/.append style={line cap=round,line join=round}}}
\pgfkeysifdefined{/tikz/alt}{}{\tikzset{alt/.code={}}}
\begin{figure}[htbp]
\centering
\begin{tikzpicture}[slfig-FIG-P684-01,>=Stealth,
  every node/.style={
    font=\fontsize{9.6pt}{11.5pt}\selectfont,
    align=center,outer sep=0pt},
  lane/.style={
    draw=SLRule,fill=SLBlueBG,rounded corners=3pt,
    line width=.78pt},
  lanehead/.style={
    draw=SLRule,fill=white,rounded corners=2.4pt,
    line width=.82pt,text width=24mm,minimum height=14mm,
    inner xsep=1.5mm,inner ysep=1.8mm,
    font=\fontsize{10.2pt}{12.2pt}\selectfont\bfseries},
  stepbox/.style={
    rounded corners=2.4pt,line width=.84pt,
    inner xsep=2.0mm,inner ysep=2.0mm,minimum height=17mm},
  hyper/.style={
    draw=SLAccentOrange,fill=SLWarmGray,dashed,
    rounded corners=2.4pt,line width=.82pt,
    text width=23mm,minimum height=14mm,
    inner xsep=1.5mm,inner ysep=1.8mm},
  randomvec/.style={
    stepbox,draw=SLMidBlue,fill=SLSoftTeal,text width=42mm},
  latentcat/.style={
    stepbox,ellipse,draw=SLDeepBlue,fill=SLSoftBlue,
    dash dot,text width=41mm},
  observed/.style={
    stepbox,rectangle,draw=SLTextGray,fill=SLWarmGray,
    double,double distance=.65pt,text width=46mm},
  variant/.style={
    draw=SLGold,fill=SLGoldBG,dashed,rounded corners=2.4pt,
    line width=.80pt,text width=32mm,
    inner xsep=1.5mm,inner ysep=1.8mm,align=left},
  prior/.style={
    -{Stealth[length=2.1mm,width=1.6mm]},
    draw=SLAccentOrange,dashed,line width=.92pt},
  main/.style={
    -{Stealth[length=2.1mm,width=1.6mm]},
    draw=SLDeepBlue,line width=.96pt},
  edgelabel/.style={
    fill=white,rounded corners=1pt,inner sep=1.1pt,
    font=\fontsize{9.6pt}{11.5pt}\selectfont},
  alt={对象—关系—结论：图保留全局主题、文档和词位三条泳道及四个生成步骤，而不复制完整 plate 图。固定超参数粗体 beta 和粗体 alpha 分别以暖色虚线框指向完整Bayes中的随机概率向量粗体 varphi_k 和粗体 theta_m；随后 theta_m 指向潜在离散变量 z_mn，z_mn 与 varphi_k 共同指向观测变量 w_mn，varphi_k 到 w_mn 的边标明按 z_mn 取一行。泳道文字说明 k 从 1 到 K、m 从 1 到 M，并且 n 从 1 到 N_m 嵌套在每个 m 内。侧注明确点参数变体中每个 varphi_k 是待估参数且不与完整Bayes混作同一后验。}]

% Three physical swimlanes.
\path[lane] (-7.65,1.35) rectangle (7.35,3.80);
\path[lane] (-7.65,-1.15) rectangle (7.35,1.15);
\path[lane] (-7.65,-4.15) rectangle (7.35,-1.40);

\node[lanehead] at (-6.20,2.58)
  {全局主题层\\$k=1,\ldots,K$};
\node[lanehead] at (-6.20,0)
  {文档层\\$m=1,\ldots,M$};
\node[lanehead,minimum height=18mm] at (-6.20,-2.78)
  {词位层\\每个 $m$ 内\\$n=1,\ldots,N_m$};

% Fixed ancestors: warm fill + dashed border + role text.
\node[hyper] (beta) at (-3.30,2.58)
  {$\boldsymbol\beta\in\mathbb R_{>0}^{V}$\\
   \textbf{固定超参数}};
\node[hyper] (alpha) at (-3.30,0)
  {$\boldsymbol\alpha\in\mathbb R_{>0}^{K}$\\
   \textbf{固定超参数}};

% Four steps.  Shape/fill/border and explicit role text remain readable in gray.
\node[randomvec] (phi) at (.60,2.58)
  {\textbf{步骤 1：抽取主题--词分布}\\
   $\boldsymbol\varphi_k\sim\operatorname{Dir}(\boldsymbol\beta)$\\
   随机概率向量\\$\boldsymbol\varphi_k\in\Delta^{V-1}$};
\node[randomvec] (theta) at (.60,0)
  {\textbf{步骤 2：抽取文档主题比例}\\
   $\boldsymbol\theta_m\sim\operatorname{Dir}(\boldsymbol\alpha)$\\
   随机概率向量\\$\boldsymbol\theta_m\in\Delta^{K-1}$};
\node[latentcat] (z) at (-1.45,-2.78)
  {\textbf{步骤 3：抽取 $z_{mn}$}\\
   $z_{mn}\sim\operatorname{Cat}(\boldsymbol\theta_m)$\\
   潜在离散变量；$z_{mn}\in\{1,\ldots,K\}$};
\node[observed] (w) at (4.05,-2.78)
  {\textbf{步骤 4：抽取 $w_{mn}$}\\
   $w_{mn}\sim\operatorname{Cat}(\boldsymbol\varphi_{z_{mn}})$\\
   观测变量；$w_{mn}\in\{1,\ldots,V\}$};

\node[variant] at (5.25,2.58)
  {\textbf{点参数变体}\\
   每个 $\boldsymbol\varphi_k$ 为待估参数；\\
   无 $\boldsymbol\beta$ 先验；\\
   非完整Bayes后验。};

% Five explicit ancestral edges.
\draw[prior] (beta.east) -- (phi.west);
\draw[prior] (alpha.east) -- (theta.west);
\draw[main] (theta.south west)
  .. controls (-.05,-1.05) and (-1.35,-1.35) .. (z.north);
\draw[main] (z.east) -- (w.west);
\draw[main] (phi.south east)
  .. controls (3.60,1.35) and (6.72,.05) ..
  node[edgelabel,pos=.62,right] {按 $z_{mn}$ 取一行}
  (w.north east);

\end{tikzpicture}
\captionsetup{width=.90\linewidth}
\caption{完整Bayes LDA 的三层祖先采样流程；点参数变体仅在侧注中区分。}
\label{fig:V5-C06-generative-process}
\end{figure}

\FloatBarrier
\noindent\textbf{读图检查。}沿
$\boldsymbol\beta\to\boldsymbol\varphi_k$、
$\boldsymbol\alpha\to\boldsymbol\theta_m$、
$\boldsymbol\theta_m\to z_{mn}$、
$z_{mn}\to w_{mn}$与$\boldsymbol\varphi_k\to w_{mn}$，核对
式\eqref{eq:V5-C06-joint}的主题词先验、文档主题先验、主题指派和发词四类因子；
再区分完整Bayes中随机的$\boldsymbol\varphi_k$与点参数变体中每个
$\boldsymbol\varphi_k$为待估参数，二者不是同一后验。

\begingroup\SetAlgoLined
% KNOWLEDGE-PLAN: KN-V5-C35-ALGORITHM_IDEA-002 L1
\paragraph{读前自检：LDA文本生成（祖先采样）：执行契约。}LDA文本生成（祖先采样）输入“文档/主题/词表规模、文档长度、严格正Dirichlet参数及可复现随机生成器状态”，条件“$M,K,V$为正整数，$N_m\in\mathbb N_0$，先验向量维数相容且有限严格为正，随机接口已登记”；请执行“概率向量完整通过单纯形检查后才原子追加”，以“完成$K+M+2I$个计划随机对象、空语料有限循环完成或首次失败时停止”判停。\par
% KNOWLEDGE-PLAN: KN-V5-C35-ALGORITHM_IDEA-003 L1
\paragraph{读前自检：LDA文本生成（祖先采样）。}LDA文本生成（祖先采样）输入“文档/主题/词表规模、文档长度、严格正Dirichlet参数及可复现随机生成器状态”，条件“$M,K,V$为正整数，$N_m\in\mathbb N_0$，先验向量维数相容且有限严格为正，随机接口已登记”；请做一轮“概率向量完整通过单纯形检查后才原子追加”，以“完成$K+M+2I$个计划随机对象、空语料有限循环完成或首次失败时停止”核对停止证书。\par

\begin{AlgorithmContract}{
  input={文档/主题/词表规模、文档长度、严格正Dirichlet参数及可复现随机生成器状态},
  output={最后合法主题/文档概率前缀与词元前缀、$k_{\rm done},m_{\rm done},i_{\rm done},q_{\rm done}$、状态与最终诊断},
  preconditions={$M,K,V$为正整数，$N_m\in\mathbb N_0$，先验向量维数相容且有限严格为正，随机接口已登记},
  initialization={初始化生成器与空输出，四个实际计数清零；总词位$I=\sum_mN_m$},
  loop={先逐主题抽$\varphi_k$，再逐文档抽$\theta_m$，最后每词位依次抽$z_{mn}$与$w_{mn}$},
  update={概率向量完整通过单纯形检查后才原子追加；一个词位的主题和词两次抽样均成功后才原子追加词元对},
  domain={$\varphi_k\in\Delta^{V-1},\theta_m\in\Delta^{K-1}$；主题/词索引合法；已提交输出有限},
  failure={维数/长度/先验非法返回$\mathtt{invalid\_input}$；随机源调用失败返回$\mathtt{random\_source\_failure}$；抽样概率、归一化或索引异常返回$\mathtt{numerical\_failure}$并保留最后合法前缀},
  stop={完成$K+M+2I$个计划随机对象、空语料有限循环完成或首次失败时停止},
  budget={至多$K$次主题Dirichlet、$M$次文档Dirichlet和$2I$次类别随机调用},
  status={$\mathtt{completed}$、$\mathtt{invalid\_input}$、$\mathtt{numerical\_failure}$、$\mathtt{random\_source\_failure}$},
  iterations={$k_{\rm done}$、$m_{\rm done}$、$i_{\rm done}$计已提交主题/文档/词位，$q_{\rm done}$计实际随机调用},
  diagnostics={随机源末状态、失败$(k)$或$(m,n,\mathrm{stage})$、概率和误差、空文档表及计划/实际工作量},
  complexity={$O(KV+MK+I(K+V))$直接类别扫描时间；保存全部输出空间$O(KV+MK+I)$}
}
\begin{algorithm}[tbp]
\caption{LDA文本生成（祖先采样）}\label{alg:V5-C06-generate}
\KwIn{$M,K,V$；长度$N_1,\ldots,N_M$；严格正的$\boldsymbol\alpha,\boldsymbol\beta$；随机数生成器状态}
\KwOut{最后合法生成前缀；$k_{\rm done},m_{\rm done},i_{\rm done},q_{\rm done}$；$\mathtt{status}$与最终诊断}
置$k_{\rm done}=m_{\rm done}=i_{\rm done}=q_{\rm done}=0$并初始化空输出；若维数、长度、先验或随机接口非法，则返回空前缀、$\mathtt{invalid\_input}$及字段诊断\;
\For{$k\leftarrow1$ \KwTo $K$}{
  调用随机源抽取临时$\varphi_k$并增加$q_{\rm done}$；调用失败时返回最后合法前缀、$\mathtt{random\_source\_failure}$及主题$k$诊断\;
  若向量非有限、非负性或归一化失败，则返回前缀、$\mathtt{numerical\_failure}$及主题概率诊断；否则原子追加并令$k_{\rm done}\leftarrow k$\;
}
\For{$m\leftarrow1$ \KwTo $M$}{
  调用随机源抽取临时$\theta_m$并增加$q_{\rm done}$；调用失败时返回前缀、$\mathtt{random\_source\_failure}$及文档$m$诊断\;
  若向量检查失败，则返回前缀、$\mathtt{numerical\_failure}$及文档概率诊断；否则原子追加并令$m_{\rm done}\leftarrow m$\;
  \For{$n\leftarrow1$ \KwTo $N_m$}{
    依次抽临时$z_{mn}$和$w_{mn}$，每次调用增加$q_{\rm done}$；任一随机调用失败则返回前缀、$\mathtt{random\_source\_failure}$及$(m,n)$阶段诊断\;
    若索引越界或类别概率检查失败，则返回前缀、$\mathtt{numerical\_failure}$及词位诊断；否则原子追加$(z_{mn},w_{mn})$并增加$i_{\rm done}$\;
  }
}
返回完整样本、实际计数、$\mathtt{completed}$及概率误差/生成器末状态最终诊断\;
\end{algorithm}
\end{AlgorithmContract}
\endgroup

\textbf{初始化与更新顺序。}初始化随机数生成器及空输出容器，先按主题编号抽取全部
$\boldsymbol\varphi_k$，再按文档编号抽取$\boldsymbol\theta_m$，最后按词位顺序先抽
$z_{mn}$、后抽$w_{mn}$。\textbf{停止条件}是所有有限循环正常完成；
\textbf{异常与退化}包括非法维数、非正先验、非有限概率及随机数生成失败，空文档是合法退化输入。
\textbf{正确性依据}是式\eqref{eq:V5-C06-joint}的祖先分解。
若保存全部输出，\textbf{空间复杂度}为$O(KV+MK+I)$；
\textbf{复杂度假设}是Dirichlet抽样与类别线性扫描都按相应向量长度计费，此时生成时间为
$O\!\left(KV+MK+I(K+V)\right)$。若先以$O(KV+MK)$时间为每个
$\boldsymbol\varphi_k$和$\boldsymbol\theta_m$建立别名表，使每次类别抽样为$O(1)$，
则总生成时间可降为$O(KV+MK+I)$。
% ALGORITHM-FIELD-NOT-APPLICABLE: alg:V5-C06-generate | 收敛条件=祖先采样只执行预先给定的有限循环，不存在迭代收敛目标; 不收敛或不适用情形=该算法生成模拟数据而不拟合观测语料，故后验不收敛问题不适用; 训练时间复杂度=文本生成不是训练任务，相关代价已单列为生成时间; 预测时间复杂度=文本生成不是给定模型后的预测接口，故预测复杂度不适用


\phantomsection\label{struct:V5-C06-CH11}
% CORE-DERIVATION-PLAN: KN-V5-C35-DERIVATION-001
\SLKnowledgeBridge{积分消去文档--主题和主题--词分布，只保留计数。}{主题指派 $z$、词元 $w$、Dirichlet超参数 $\alpha,\beta$ 及计数 $n_{mk},n_{kv}$。}{$\alpha_k,\beta_v>0$；删去当前词元后的计数用上标 $-i$；采用多元 Beta函数。}{先按文档和主题分组联合分布，并分别应用 \DirichletMultinomial{}积分。}

\begin{derivationgoalbox}
从式\eqref{eq:V5-C06-joint}积分消去
$\boldsymbol\theta$和$\boldsymbol\varphi$，得到只依赖计数的
$p(\boldsymbol w,\boldsymbol z\mid\boldsymbol\alpha,\boldsymbol\beta)$；
随后用“加入一个词元前后”的比值推出Gibbs满条件分布。
\end{derivationgoalbox}
\begin{derivationprepbox}
记$\alpha_0=\sum_k\alpha_k$、$\beta_0=\sum_v\beta_v$；
$n_{mk}=\sum_n\mathbf1\{z_{mn}=k\}$，
$n_{kv}=\sum_{m,n}\mathbf1\{z_{mn}=k,w_{mn}=v\}$，
$n_{m\cdot}=\sum_kn_{mk}$，$n_{k\cdot}=\sum_vn_{kv}$；并记
$\boldsymbol n_m=(n_{m1},\ldots,n_{mK})$、
$\boldsymbol n_k=(n_{k1},\ldots,n_{kV})$。上标$^{-i}$表示删去当前词元$i=(m,n)$后的计数。
多元Beta函数记为$B(\boldsymbol a)=\prod_j\Gamma(a_j)/\Gamma(\sum_ja_j)$。
\end{derivationprepbox}

\textbf{推导第1步：按主题收集单词似然。}
\[
p(\boldsymbol w\mid\boldsymbol z,\boldsymbol\varphi)
=\prod_{k=1}^{K}\prod_{v=1}^{V}\varphi_{kv}^{n_{kv}}.
\]
由定义\ref{def:V5-C06-lda}中的$\beta_v>0$，每个主题的先验密度与有序词元似然的乘积为
\[
\frac1{B(\boldsymbol\beta)}
\prod_{v=1}^V\varphi_{kv}^{\beta_v-1}
\prod_{v=1}^V\varphi_{kv}^{n_{kv}}
=\frac1{B(\boldsymbol\beta)}
\prod_{v=1}^V\varphi_{kv}^{n_{kv}+\beta_v-1}.
\]
在$\Delta_{V-1}$上逐项积分，并使用定理\ref{thm:V5-C05-conjugacy}背后的多元Beta恒等式，得到
\[
\int_{\Delta_{V-1}}
\frac1{B(\boldsymbol\beta)}
\prod_{v=1}^V\varphi_{kv}^{n_{kv}+\beta_v-1}
\,\mathrm d\boldsymbol\varphi_k
=\frac{B(\boldsymbol n_k+\boldsymbol\beta)}{B(\boldsymbol\beta)}.
\]
这里$\boldsymbol w,\boldsymbol z$表示按位置编号的词元与主题指派，故似然是逐位置概率的乘积，没有计数向量概率中的多项组合系数。对$k=1,\ldots,K$相乘得到
\begin{equation}
p(\boldsymbol w\mid\boldsymbol z,\boldsymbol\beta)
=\prod_{k=1}^{K}
\frac{B(\boldsymbol n_k+\boldsymbol\beta)}{B(\boldsymbol\beta)}.
\label{eq:V5-C06-collapsed-word}
\end{equation}
\textbf{推导第2步：按文档收集主题指派。}把上一行中的主题索引$k$与词索引$v$分别换成文档索引$m$与主题索引$k$，并由$\alpha_k>0$，对每篇文档逐项写出
\[
\int_{\Delta_{K-1}}
\frac1{B(\boldsymbol\alpha)}
\prod_{k=1}^K\theta_{mk}^{n_{mk}+\alpha_k-1}
\,\mathrm d\boldsymbol\theta_m
=\frac{B(\boldsymbol n_m+\boldsymbol\alpha)}{B(\boldsymbol\alpha)}.
\]
这仍是有序主题指派的逐位置似然积分；没有额外的多项组合系数。对$m=1,\ldots,M$相乘，积分消去$\boldsymbol\theta_m$：
\begin{equation}
p(\boldsymbol z\mid\boldsymbol\alpha)
=\prod_{m=1}^{M}
\frac{B(\boldsymbol n_m+\boldsymbol\alpha)}{B(\boldsymbol\alpha)}.
\label{eq:V5-C06-collapsed-topic}
\end{equation}

\phantomsection\label{struct:V5-C06-CH12}
\begin{theorem}[LDA的折叠联合分布]\label{thm:V5-C06-collapsed-joint}
在定义\ref{def:V5-C06-lda}的条件下，
\begin{equation}
p(\boldsymbol w,\boldsymbol z\mid\boldsymbol\alpha,\boldsymbol\beta)
=\prod_{k=1}^{K}
\frac{B(\boldsymbol n_k+\boldsymbol\beta)}{B(\boldsymbol\beta)}
\prod_{m=1}^{M}
\frac{B(\boldsymbol n_m+\boldsymbol\alpha)}{B(\boldsymbol\alpha)}.
\label{eq:V5-C06-collapsed-joint}
\end{equation}
该式保留了与$\boldsymbol z$有关的全部归一化常数。
\end{theorem}
% KNOWLEDGE-PLAN: KN-V5-C35-PROOF-002 L1
\paragraph{读前自检：证明：LDA的折叠联合分布。}LDA折叠联合分布的被积函数非负，且主题参数与文档参数两组积分变量互不重叠；请分别写出逐主题、逐文档的多元Beta比值并相乘，核对$\alpha_k,\beta_v>0$保证积分有限且无需多项组合系数。\par

\begin{proof}
从式\eqref{eq:V5-C06-joint}出发，被积函数非负，Tonelli定理允许改变积分次序。上面的逐主题计算给出式\eqref{eq:V5-C06-collapsed-word}，逐文档计算给出式\eqref{eq:V5-C06-collapsed-topic}。两组积分变量互不重叠，故两式相乘即得\eqref{eq:V5-C06-collapsed-joint}。所有$\alpha_k,\beta_v$严格为正，确保每个多元Beta积分有限；由于对象是有序词元与指派的联合概率，整个推导中都不出现多项组合系数。
\end{proof}

\phantomsection\label{struct:V5-C06-CH13}
\begin{corollary}[文档内主题指派的交换性]\label{cor:V5-C06-topic-exchangeable}
固定观测词和其他文档后，若同时置换文档$m$的词位置及相应主题指派，式
\eqref{eq:V5-C06-collapsed-joint}不变。
\end{corollary}
% KNOWLEDGE-PLAN: KN-V5-C35-PROOF-003 L1
\paragraph{读前自检：证明：文档内主题指派的交换性。}LDA折叠联合分布只通过$n_{mk}$与$n_{kv}$依赖位置数据；请同步置换词元和主题指派，核对两组计数及Beta比值不变，再只置换词元构造$n_{kv}$可能改变的对照。\par

\begin{proof}
折叠联合分布只通过$n_{mk}$和$n_{kv}$依赖位置数据；同步置换不改变任一计数，
故每个Beta函数比值不变。若只置换词而不随之置换指派，$n_{kv}$可能改变，结论不成立。
\end{proof}

% KNOWLEDGE-PLAN: KN-V5-C35-ALGORITHM_IDEA-004 L1
\paragraph{读前自检：LDA折叠Gibbs抽样。}LDA折叠Gibbs对词元$i=(m,n)$、观测词$v$采用留一计数；请对每个主题计算$(n_{mk}^{-i}+\alpha_k)(n_{kv}^{-i}+\beta_v)/(n_{k\cdot}^{-i}+\beta_0)$，再对$k$归一化，并核对得到合法满条件概率。\par
% KNOWLEDGE-PLAN: KN-V5-C35-CONCEPT-007 L1
\paragraph{读前自检：折叠抽样的留一计数思想。}折叠抽样更新词元$i$前必须先从文档--主题和主题--词计数表减去其旧主题；请执行一次“减一”，核对两张表总计数各减少1，再用留一计数计算候选权重。\par
% KNOWLEDGE-PLAN: KN-V5-C35-BOUNDARY-001 L1
\paragraph{读前自检：LDA折叠Gibbs条件更新。}LDA折叠Gibbs满条件由文档主题计数因子与主题词计数因子相乘；请为全部$k$算未归一化权重并除以权重和，核对概率非负且总和为1。\par
% KNOWLEDGE-PLAN: KN-V5-C35-FORMULA-002 L1
\paragraph{读前自检：主题分布与主题词分布的后验汇总。}折叠Gibbs样本的后验汇总用$\widehat\theta_{mk}=(n_{mk}+\alpha_k)/(N_m+\alpha_0)$和$\widehat\varphi_{kv}=(n_{kv}+\beta_v)/(n_{k\cdot}+\beta_0)$；请逐行求和核对为1。\par
% KNOWLEDGE-PLAN: KN-V5-C35-ALGORITHM_IDEA-005 L1
\paragraph{读前自检：折叠Gibbs扫描流程。}折叠Gibbs扫描以当前主题指派和一致计数表为状态；请对一个词元完成“减--算--抽--加”，核对两张计数表与新指派一致，整轮提交后按扫描预算停止并保留链诊断。\par
% KNOWLEDGE-PLAN: KN-V5-C35-CONCEPT-008 L1
\paragraph{读前自检：主题指派与计数表。}主题指派$z_i$与计数$n_{mk},n_{kv},n_{k\cdot}$必须同步；请改变一个词元的主题，逐项更新受影响的旧、新主题计数，核对文档总数和主题总数守恒。\par

\SLTeachSection{折叠Gibbs推断}{sec:V5-C06-S03}
\knowledgeanchor{SLM-20-03-001}{LDA折叠Gibbs抽样}
\knowledgeanchor{SLM-20-03-002}{折叠抽样的留一计数思想}
\knowledgeanchor{SLM-20-03-003}{LDA折叠Gibbs条件更新}
\knowledgeanchor{SLM-20-03-004}{主题分布与主题词分布的后验汇总}
\knowledgeanchor{SLM-20-03-005}{折叠Gibbs扫描流程}
\knowledgeanchor{SLM-20-02-008}{主题指派与计数表}

\begin{proposition}[折叠Gibbs满条件分布]\label{prop:V5-C06-gibbs-conditional}
设当前词元$i=(m,n)$的观测词为$v$。删去其旧主题后的满条件概率满足
\begin{equation}
p(z_i=k\mid\boldsymbol z_{-i},\boldsymbol w,
\boldsymbol\alpha,\boldsymbol\beta)
\propto
\frac{n_{kv}^{-i}+\beta_v}{n_{k\cdot}^{-i}+\beta_0}
\frac{n_{mk}^{-i}+\alpha_k}{n_{m\cdot}^{-i}+\alpha_0}.
\label{eq:V5-C06-gibbs-conditional}
\end{equation}
右端对$k=1,\ldots,K$归一化后才是概率。
\end{proposition}
% KNOWLEDGE-PLAN: KN-V5-C35-PROOF-004 L1
\paragraph{读前自检：证明：折叠Gibbs满条件分布。}折叠Gibbs比较候选主题$k$时只保留主题--词与文档--主题计数因子；请用$\Gamma(x+1)/\Gamma(x)=x$分别约去两个Beta函数比值，核对所得两项相乘即为未归一化满条件权重。\par

\begin{proof}
在式\eqref{eq:V5-C06-collapsed-joint}中，除主题$k$的词计数因子和文档$m$的主题计数因子外，其余因子在比较候选$k$时相消。
使用$\Gamma(x+1)/\Gamma(x)=x$，把“加入当前词元”后的Beta函数除以留一Beta函数，分别得到
\[
\frac{n_{kv}^{-i}+\beta_v}{n_{k\cdot}^{-i}+\beta_0},
\qquad
\frac{n_{mk}^{-i}+\alpha_k}{n_{m\cdot}^{-i}+\alpha_0}.
\]
两项相乘即为未归一化权重。第二项分母虽对$k$相同，保留它能显示完整的后验预测来源。
\end{proof}

\textbf{推导第3步：用留一状态的Beta函数比值形成满条件。}
命题\ref{prop:V5-C06-gibbs-conditional}把折叠联合分布中所有与候选主题无关的因子约去，
再对$K$个正权重统一归一化；这一“先留一、后作比”的顺序也是实现计数守恒的依据。

图\ref{fig:V5-C06-collapsed-gibbs-counts}把一次更新需要读写的两张计数表并列画出：
先从旧主题的文档--主题计数和主题--词计数各减一，再计算所有候选权重，抽取新主题后再各加一。
% v2.6.0 figure UID: FIG-P687-01
% Iteration 03: widen the badge-2-to-document-card whitespace corridor.
\tikzset{slfig-FIG-P687-01/.style={font=\fontsize{9.2pt}{11.0pt}\selectfont,every path/.append style={line cap=round,line join=round}}}
\pgfkeysifdefined{/tikz/alt}{}{\tikzset{alt/.code={}}}
\begin{figure}[htbp]
\centering
\begin{tikzpicture}[slfig-FIG-P687-01,
  >=Stealth,
  every node/.style={font=\fontsize{9.2pt}{11pt}\selectfont,align=center,outer sep=0pt},
  action/.style={draw=SLRuleGray,fill=SLSoftGray,rounded corners=2pt,
    minimum height=12mm,text width=36mm,inner xsep=3mm,inner ysep=2mm},
  evidence/.style={draw=SLMidBlue,fill=SLSoftBlue,rounded corners=2pt,
    minimum height=16mm,text width=44mm,inner xsep=3mm,inner ysep=2mm},
  conditional/.style={draw=SLDeepBlue,fill=white,rounded corners=2pt,
    minimum height=17mm,text width=61mm,inner xsep=3mm,inner ysep=2mm},
  badge/.style={circle,fill=SLDeepBlue,text=white,minimum size=6.4mm,inner sep=0pt,
    font=\fontsize{9.0pt}{10.6pt}\selectfont\bfseries},
  branch/.style={-{Stealth[length=2.1mm,width=1.55mm]},draw=SLMidBlue,line width=.78pt},
  loop/.style={-{Stealth[length=2.2mm,width=1.6mm]},draw=SLDeepBlue,line width=.95pt},
  alt={对象—关系—结论：折叠Gibbs更新构成五步闭环：先移除当前词位，再分别读取带负i上标的文档主题计数和主题词计数，两支证据相乘归一化得到满条件，最后抽样并恢复计数后进入下一词位。}]

\node[action] (remove) at (0,3.00)
  {移除当前词位$i$\\从两张计数表各减一};
\node[badge] at (-2.78,3.50) {1};

\node[evidence] (doc) at (-3.75,1.05)
  {读取文档--主题计数\\
   $\displaystyle n_{mk}^{-i}+\alpha_k$\\[-1pt]
   $\displaystyle n_{m\cdot}^{-i}+\alpha_0$};
\node[badge] at (-6.75,1.48) {2};
\node[evidence] (word) at (3.75,1.05)
  {读取主题--单词计数\\
   $\displaystyle n_{kv}^{-i}+\beta_v$\\[-1pt]
   $\displaystyle n_{k\cdot}^{-i}+\beta_0$};
\node[badge] at (0.68,1.48) {3};

\node[conditional] (full) at (0,-1.00)
  {相乘并对$k$归一化\\[-1pt]
   $\displaystyle p(z_i=k\mid-)
   \propto
   \frac{n_{mk}^{-i}+\alpha_k}{n_{m\cdot}^{-i}+\alpha_0}
   \frac{n_{kv}^{-i}+\beta_v}{n_{k\cdot}^{-i}+\beta_0}$};
\node[badge] at (-4.05,-0.53) {4};

\node[action,draw=SLDeepBlue,fill=SLSoftBlue,text width=41mm] (sample) at (0,-2.85)
  {抽样$k^*$并恢复计数\\把词位$i$加回两张表};
\node[badge] at (-3.03,-2.43) {5};

\draw[branch] (remove.south west) -- (doc.north);
\draw[branch] (remove.south east) -- (word.north);
\draw[branch] (doc.south) -- (full.north west);
\draw[branch] (word.south) -- (full.north east);
\draw[loop] (full.south) -- (sample.north);
\draw[loop] (sample.east) -- (6.65,-2.85) --
  node[midway,right=2.4mm,text width=20mm,align=left,font=\fontsize{8.8pt}{10.5pt}\selectfont,text=SLTextGray]
    {下一词位：\\重新从“移除”开始}
  (6.65,3.00) -- (remove.east);

\node[font=\fontsize{9.0pt}{10.6pt}\selectfont,text=SLTextGray] at (0,-3.90)
  {上标$-i$始终保留到抽样完成，防止当前词位对自身重复计数。};
\end{tikzpicture}
\caption{折叠Gibbs更新先临时移除当前词位，再把文档--主题偏好与主题--单词证据相乘形成满条件分布，最后抽取新主题并恢复计数；上标负$i$防止当前词位对自身产生重复计数}
\label{fig:V5-C06-collapsed-gibbs-counts}
\end{figure}


样本之后可用后验均值平滑估计
\begin{align}
\widehat\theta_{mk}
&=\frac{n_{mk}+\alpha_k}{n_{m\cdot}+\alpha_0},
&
\widehat\varphi_{kv}
&=\frac{n_{kv}+\beta_v}{n_{k\cdot}+\beta_0}.
\label{eq:V5-C06-gibbs-postprocess}
\end{align}
这些是给定某个保留样本的后验均值；更稳健的做法是对多个保留样本的预测量或计数平滑估计取平均，而不是只报告最后一次扫描。

\begin{implementationbox}
实现折叠Gibbs抽样时，输入应同时包括$\boldsymbol w$、$\boldsymbol\alpha$、
$\boldsymbol\beta$与$K$，避免把模型参数与数据割裂。
燃烧期结束后还要继续留样，
可按预先规定或诊断选择间隔，并报告有效样本量、链间差异和随机种子。燃烧期不是算法完成的同义词。
\end{implementationbox}

% KNOWLEDGE-PLAN: KN-V5-C35-ALGORITHM_IDEA-006 L1
\paragraph{读前自检：LDA变分EM局部—全局更新。}LDA变分EM在局部$q_m(z_m,\theta_m)$与全局点参数$\varphi,\alpha$间交替；请完成一次局部责任度归一化，再用期望计数更新一行$\varphi_k$，核对ELBO不下降后提交。\par
% KNOWLEDGE-PLAN: KN-V5-C35-ALGORITHM_IDEA-007 L1
\paragraph{读前自检：平均场变分推断。}平均场变分族把隐变量分解为若干$q_j(h_j)$；请固定其余因子并计算$q_j^*(h_j)\propto\exp(E_{-j}\log p(w,h))$，核对归一化常数有限。\par
% KNOWLEDGE-PLAN: KN-V5-C35-ALGORITHM_IDEA-008 L1
\paragraph{读前自检：LDA变分EM算法。}LDA变分EM以语料、$K$、正$\alpha$与列随机$\varphi$为输入；请完成一篇文档的局部坐标迭代和一次全局期望计数更新，只在ELBO/归一化门禁通过时提交，按双容差或预算停止。\par
% KNOWLEDGE-PLAN: KN-V5-C35-ALGORITHM_IDEA-009 L1
\paragraph{读前自检：通用变分EM坐标上升。}通用变分EM坐标上升固定其余$q_{-j}$时更新$q_j$，再固定$q$更新模型参数；请执行一个坐标最优更新，核对归一化且ELBO不降后再进入M步。\par
% KNOWLEDGE-PLAN: KN-V5-C35-CONCEPT-009 L1
\paragraph{读前自检：LDA变分更新推导。}LDA变分更新从$\log p(w,h)$中保留含当前因子的项并对其余变量取期望；请指数化并归一化一次，核对得到的责任度非负且和为1。\par
% KNOWLEDGE-PLAN: KN-V5-C35-ALGORITHM_IDEA-010 L1
\paragraph{读前自检：LDA的变分参数估计算法。}LDA变分参数估计维护局部责任度/Dirichlet参数与全局主题词分布；请用局部期望计数更新一个$\varphi_{kv}$并归一化，核对ELBO不降，按外层容差或预算停止。\par

\SLAdvancedSection{变分EM推断}{sec:V5-C06-S04}
\knowledgeanchor{SLM-20-04-001}{LDA变分EM局部—全局更新}
\knowledgeanchor{SLM-20-04-002}{平均场变分推断}
\knowledgeanchor{SLM-20-04-003}{LDA变分EM算法}
\knowledgeanchor{SLM-20-03-006}{通用变分EM坐标上升}
\knowledgeanchor{SLM-20-04-004}{LDA变分更新推导}
\knowledgeanchor{SLM-20-04-005}{LDA的变分参数估计算法}

\begin{theorem}[KL--ELBO恒等式]\label{thm:V5-C06-elbo}
设$0<p(\boldsymbol w)<\infty$，联合密度与$q$相对于同一支配测度可测，$q\ll p(\boldsymbol h\mid\boldsymbol w)$，并且
\[
\mathbb E_q\!\left|\log p(\boldsymbol w,\boldsymbol h)\right|<\infty,
\qquad
\mathbb E_q\!\left|\log q(\boldsymbol h)\right|<\infty.
\]
这些条件排除了未定义的$\infty-\infty$。则
\begin{equation}
\log p(\boldsymbol w)
=\mathcal L(q)+
\KL\!\left(q(\boldsymbol h)\,
\middle\|\,p(\boldsymbol h\mid\boldsymbol w)\right),
\label{eq:V5-C06-elbo-identity}
\end{equation}
其中$\boldsymbol h$汇总隐变量，且
\[
\mathcal L(q)=
\mathbb E_q\log p(\boldsymbol w,\boldsymbol h)
-\mathbb E_q\log q(\boldsymbol h).
\]
因此$\mathcal L(q)\le\log p(\boldsymbol w)$，等号当且仅当$q$等于精确后验（几乎处处）。
\end{theorem}
% KNOWLEDGE-PLAN: KN-V5-C35-PROOF-005 L1
\paragraph{读前自检：证明：KL--ELBO恒等式。}KL--ELBO恒等式假定$0<p(\boldsymbol w)<\infty$、$q$对后验绝对连续且相关项绝对可积；请把$\log p(\boldsymbol h\mid\boldsymbol w)=\log p(\boldsymbol w,\boldsymbol h)-\log p(\boldsymbol w)$代入KL，核对$\KL(q\|p)=\log p(\boldsymbol w)-\mathcal L(q)$。\par

\begin{proof}
由$0<p(\boldsymbol w)<\infty$与绝对连续性，Bayes等式在$q$几乎处处成立：
$\log p(\boldsymbol h\mid\boldsymbol w)
=\log p(\boldsymbol w,\boldsymbol h)-\log p(\boldsymbol w)$。
上述绝对可积条件允许分别积分并整理，不会出现$\infty-\infty$：
\begin{align*}
\KL(q\|p)
&=\mathbb E_q\log q(\boldsymbol h)
-\mathbb E_q\log p(\boldsymbol h\mid\boldsymbol w)\\
&=\log p(\boldsymbol w)-\mathcal L(q).
\end{align*}
KL非负给出下界与等号条件。
\end{proof}

图\ref{fig:V5-C06-elbo-geometry}显示：受限族$\mathcal Q$中最优$q^*$通常仍与真实后验有正KL距离，故“E步后ELBO等于证据”一般错误。
% v2.6 layout: former forced clear removed; natural pagination.
% v2.6.0 figure UID: FIG-P689-01
% Iteration 05: use an explicit 8 mm x-shift to move the complete right-hand axis into its panel.
\tikzset{slfig-FIG-P689-01/.style={font=\fontsize{9.2pt}{11.0pt}\selectfont,every path/.append style={line cap=round,line join=round}}}
\pgfplotsset{slfig-FIG-P689-01-axis/.style={tick label style={font=\footnotesize},label style={font=\small},clip=false}}
\pgfkeysifdefined{/tikz/alt}{}{\tikzset{alt/.code={}}}
\begin{figure}[htbp]
\centering
\begin{tikzpicture}[slfig-FIG-P689-01,
  every node/.style={font=\fontsize{9.2pt}{11.0pt}\selectfont,align=center,outer sep=0pt},
  panel/.style={draw=SLRuleGray,line width=.55pt,rounded corners=3pt,
    minimum width=67mm,minimum height=53mm,inner sep=0pt},
  alt={对象—关系—结论：左侧以长度分解直接显示观测对数证据等于ELBO与非负KL间隙之和；右侧坐标更新使ELBO形成非降阶梯，但只能到坐标稳定或局部驻点，未知全局上限不应被解释为已经达到。}]

\node[panel] at (-3.85,0.10) {};
\node[font=\bfseries,text=SLDeepBlue] at (-3.85,2.42)
  {证据的长度分解};
\fill[SLMidBlue,opacity=0.08] (-6.55,0.50) rectangle (-2.55,1.25);
\fill[fill=SLAccentOrange,opacity=0.08] (-2.55,0.50) rectangle (-0.95,1.25);
\draw[draw=SLDeepBlue,line width=.95pt] (-6.55,0.50) rectangle (-0.95,1.25);
\draw[draw=SLMidBlue,line width=.90pt] (-2.55,0.50) -- (-2.55,1.25);
\node[text=SLDeepBlue] at (-4.55,0.875) {$\mathcal L(q)$：证据下界};
\node[text=SLAccentOrange] at (-1.75,0.875) {KL\ 间隙};
\draw[<->,draw=SLTextGray,line width=.58pt] (-6.55,1.46) -- (-0.95,1.46)
  node[midway,above=2.0mm] {$\log p(w)$};
\node[align=left,text width=59mm] at (-3.85,-0.15)
  {$\displaystyle \log p(w)=\mathcal L(q)+
    \operatorname{KL}\!\left(q(h)\,\|\,p(h\mid w)\right)$};
\node[align=left,text width=59mm,text=SLTextGray] at (-3.85,-1.35)
  {KL$\ge0$，故$\mathcal L(q)\le\log p(w)$；变分族受限时，间隙可保持为正。};

\node[panel] at (3.85,0.10) {};
\node[font=\bfseries,text=SLDeepBlue] at (3.85,2.42)
  {坐标更新下的 ELBO 非降阶梯};
\begin{axis}[slfig-FIG-P689-01-axis,
  at={(0.80,-1.55)},anchor=south west,xshift=8mm,yshift=-10.5mm,
  width=59mm,height=34mm,scale only axis,
  xmin=0,xmax=6.3,ymin=0,ymax=1.12,
  axis lines=left,axis line style={draw=SLTextGray,line width=.55pt},
  tick style={draw=SLRuleGray},xtick={0,1,2,3,4,5,6},ytick=\empty,
  xlabel={坐标更新轮次},xlabel style={font=\fontsize{9.0pt}{10.6pt}\selectfont},
  clip=false]
  \addplot+[const plot,draw=SLMidBlue,line width=.95pt,mark=*,mark size=1.5pt]
    coordinates {(0,0.12) (1,0.34) (2,0.50) (3,0.64) (4,0.73) (5,0.79) (6,0.80)};
  \addplot+[draw=SLAccentOrange,dashed,line width=.70pt,mark=none]
    coordinates {(0,1.02) (6.3,1.02)};
  \node[anchor=north west,font=\fontsize{9.0pt}{10.6pt}\selectfont,text=SLAccentOrange]
    at (axis cs:0.15,0.98) {未知全局上限};
  \node[anchor=north east,font=\fontsize{9.0pt}{10.6pt}\selectfont,text=SLDeepBlue]
    at (axis cs:6.18,0.48) {坐标稳定／局部驻点};
\end{axis}
\end{tikzpicture}
\captionsetup{width=.94\linewidth}
\caption{观测对数证据等于ELBO与变分分布到真实后验的KL散度之和；平均场坐标更新可使ELBO逐步不降，但非凸目标下有限运行通常只得到坐标稳定点或局部驻点，多启动比较也不构成全局最优证明}
\label{fig:V5-C06-elbo-geometry}
\end{figure}


以下变分推导只针对式\eqref{eq:V5-C06-parameterized-joint}的点参数变体；
$\boldsymbol\beta$不属于该变体的ELBO。为使该模型的对数项、坐标命题与评价口径使用同一支持，先冻结训练活动词表
$\mathcal V_{\mathrm{tr}}$：每个$v\in\mathcal V_{\mathrm{tr}}$在训练语料中至少出现一次，测试时的词表外项要么映射到一个已在训练中出现的显式
$\mathtt{UNK}$，要么返回$\mathtt{invalid\_input}(\mathtt{support\_failure},v)$。本点参数模型的可行域规定
\[
\varphi_{kv}>0\quad(k=1,\ldots,K,\ v\in\mathcal V_{\mathrm{tr}}),
\qquad \sum_{v\in\mathcal V_{\mathrm{tr}}}\varphi_{kv}=1.
\]
这是模型支持的明示限制，不是在实现中偷偷加入$\boldsymbol\beta$平滑。对固定$\boldsymbol\varphi$的单文档，采用平均场族
\begin{equation}
q(\boldsymbol\theta_m,\boldsymbol z_m
\mid\boldsymbol\gamma_m,\boldsymbol\eta_m)
=\operatorname{Dir}(\boldsymbol\theta_m\mid\boldsymbol\gamma_m)
\prod_{n=1}^{N_m}\operatorname{Cat}(z_{mn}\mid\boldsymbol\eta_{mn}).
\label{eq:V5-C06-mean-field}
\end{equation}
图\ref{fig:V5-C06-mean-field-graph}强调：近似族主动切断了后验中
$\boldsymbol\theta_m$与各$z_{mn}$的依赖，这换来可计算性，也造成近似误差。
% v2.6 layout: former forced clear removed; natural pagination.
% v2.6.0 figure UID: FIG-P690-01
% Mean-field comparison: real formula bands, an unobstructed transition corridor,
% and a geometric break instead of an opaque white eraser.
\tikzset{slfig-FIG-P690-01/.style={font=\fontsize{9.2pt}{11.0pt}\selectfont,every path/.append style={line cap=round,line join=round}}}
\pgfkeysifdefined{/tikz/alt}{}{\tikzset{alt/.code={}}}
\begin{figure}[htbp]
\centering
\begin{tikzpicture}[slfig-FIG-P690-01,
  >=Stealth,
  every node/.style={font=\fontsize{9.2pt}{11.0pt}\selectfont,align=center},
  latent/.style={circle,draw=SLMidBlue,fill=white,minimum size=10mm,inner sep=1pt},
  qlatent/.style={latent,minimum size=15.5mm,font=\fontsize{9.2pt}{11.0pt}\selectfont},
  observed/.style={circle,draw=SLDeepBlue,fill=SLDeepBlue,text=white,minimum size=10mm,inner sep=1pt},
  param/.style={draw=SLRuleGray,fill=SLSoftGray,rounded corners=2pt,minimum width=20mm,minimum height=8mm},
  panel/.style={draw=SLRuleGray,rounded corners=3pt,densely dashed,inner xsep=5.8mm,inner ysep=5.5mm},
  dep/.style={draw=SLMidBlue,line width=0.90pt},
  input/.style={->,draw=SLDeepBlue,line width=0.80pt},
  cutedge/.style={draw=SLRuleGray,densely dashed,line width=0.55pt},
  cutmark/.style={draw=SLAccentOrange,line width=0.70pt},
  >={Stealth[length=2.2mm,width=1.6mm]},
  alt={对象—关系—结论：左右面板使用完全相同的节点位置。真实单文档后验中主题比例与主题指派耦合；平均场近似切断这一后验依赖，但仍保留观测词与固定主题词参数对责任度更新的输入。}]

\node[font=\bfseries\fontsize{9.6pt}{11.5pt}\selectfont,text=SLDeepBlue] at (-4.40,2.55)
  {真实后验：$\theta_m$ 与 $z_{mn}$ 耦合};
\node[latent] (tl) at (-6.40,0.75) {$\theta_m$};
\node[latent] (zl) at (-4.40,0.75) {$z_{mn}$};
\node[observed] (wl) at (-2.40,0.75) {$w_{mn}$};
\node[param] (pl) at (-4.40,-0.85) {$\boldsymbol\varphi$ 固定};
\draw[dep] (tl) -- (zl);
\draw[input] (wl) -- (zl);
\draw[input] (pl) -- (zl);
\node[panel,fit=(tl)(zl)(wl)(pl)] (leftpanel) {};
\node[text=SLTextGray] at (-4.40,-2.25)
  {$p(\theta_m,z_{mn}\mid w_m,\boldsymbol\varphi)$};

\node[font=\bfseries\fontsize{9.6pt}{11.5pt}\selectfont,text=SLDeepBlue] at (4.40,2.55)
  {平均场：切断近似后验依赖};
\node[qlatent] (tr) at (2.40,0.75) {$q(\theta_m)$};
\node[qlatent] (zr) at (4.40,0.75) {$q(z_{mn})$};
\node[observed] (wr) at (6.40,0.75) {$w_{mn}$};
\node[param] (pr) at (4.40,-0.85) {$\boldsymbol\varphi$ 固定};
\draw[cutedge] (tr.east) -- (3.21,0.75);
\draw[cutedge] (3.59,0.75) -- (zr.west);
\draw[cutmark] (3.34,0.59) -- (3.39,0.89);
\draw[cutmark] (3.41,0.59) -- (3.46,0.89);
\draw[input] (wr) -- (zr);
\draw[input] (pr) -- (zr);
\node[panel,fit=(tr)(zr)(wr)(pr)] (rightpanel) {};
\node[text=SLTextGray] at (4.40,-2.25)
  {$q(\theta_m)\prod_n q(z_{mn})$};

\draw[->,draw=SLDeepBlue,line width=0.85pt] (leftpanel.east) -- (rightpanel.west)
  node[midway,above=2.2mm,font=\fontsize{9.2pt}{10.8pt}\selectfont,
       text width=26mm,align=center]{选择可计算的\\乘积分布族};
\node[draw=SLMidBlue,fill=SLSoftBlue,rounded corners=2pt,
  text width=125mm,minimum height=10mm,inner xsep=3.0mm,inner ysep=1.8mm] at (0,-3.84)
  {切断的是近似后验中的直接依赖，不是删除生成模型依赖；\\
   $w_{mn}$ 与$\boldsymbol\varphi$仍进入$q(z_{mn})$的责任度更新。};
\end{tikzpicture}
\caption{固定主题--词参数$\boldsymbol\varphi$时，平均场近似把单文档后验中耦合的$\boldsymbol\theta_m$与$z_{mn}$替换为两个变分因子族；被切断的是近似后验中的直接依赖，固定参数仍通过词似然进入责任度更新}
\label{fig:V5-C06-mean-field-graph}
\end{figure}


% KNOWLEDGE-PLAN: KN-V5-C35-PROPOSITION-003 L1
\paragraph{读前自检：平均场坐标最优形式。}平均场坐标最优形式以“设$A$关于支配测度$\mu_j$几乎处处为实值且可测”为约束；请把$q_j^*(h_j)$展开一次并检查其类型或边界，并确认展开后的量满足定义中的域、维数或符号限制。\par

\begin{proposition}[平均场坐标最优形式]\label{prop:V5-C06-mean-field-coordinate}
固定其他变分因子，令$A(h_j)=\mathbb E_{q_{-j}}\log p(\boldsymbol w,\boldsymbol h)$。设$A$关于支配测度$\mu_j$几乎处处为实值且可测，并且
\[
0<Z_j=\int e^{A(h_j)}\,\mu_j(\mathrm dh_j)<\infty.
\]
再假定
\[
\int e^{A(h_j)}\lvert A(h_j)\rvert\,\mu_j(\mathrm dh_j)<\infty.
\]
则该坐标的正规化最优解存在并满足
\begin{equation}
q_j^*(h_j)=\frac{e^{A(h_j)}}{Z_j},
\qquad
\log q_j^*(h_j)=A(h_j)-\log Z_j,
\label{eq:V5-C06-coordinate-general}
\end{equation}
并在$\mu_j$零测集意义下唯一。
\end{proposition}
% KNOWLEDGE-PLAN: KN-V5-C35-PROOF-006 L1
\paragraph{读前自检：证明：平均场坐标最优形式。}平均场坐标更新只比较$q_j\ll q_j^*$且熵项、$A(h_j)$项可积的密度；请将坐标目标改写为$\log Z_j-\KL(q_j\|q_j^*)$，核对$q_j=q_j^*$几乎处处时唯一取到最大值。\par

\begin{proof}
把ELBO中与$q_j$有关的部分记为
\[
\mathcal F(q_j)=\int q_j(h_j)A(h_j)\,\mu_j(\mathrm dh_j)
-\int q_j(h_j)\log q_j(h_j)\,\mu_j(\mathrm dh_j),
\]
并只比较满足$q_j\ll q_j^*$、$\int q_j\lvert A\rvert\,\mathrm d\mu_j<\infty$且
$\int q_j\lvert\log q_j\rvert\,\mathrm d\mu_j<\infty$的概率密度。新增的加权可积条件给出
\[
\int q_j^*\lvert A\rvert\,\mathrm d\mu_j
=Z_j^{-1}\int e^A\lvert A\rvert\,\mathrm d\mu_j<\infty,
\qquad
\int q_j^*\lvert\log q_j^*\rvert\,\mathrm d\mu_j<\infty,
\]
所以$q_j^*$本身属于比较类，$\mathcal F(q_j^*)$不是$\infty-\infty$。由$\log q_j^*=A-\log Z_j$，逐项代入可得
\[
\mathcal F(q_j^*)-\mathcal F(q_j)
=\int q_j\log\frac{q_j}{q_j^*}\,\mathrm d\mu_j
=\KL(q_j\|q_j^*)\ge0.
\]
因此$q_j^*$达到最大值；等号当且仅当$q_j=q_j^*$几乎处处，故最优解唯一。若$Z_j=0$或$Z_j=\infty$，则此正规化坐标更新不存在；即使$0<Z_j<\infty$，若$\int e^A\lvert A\rvert\,\mathrm d\mu_j=\infty$，上述两项分开定义的ELBO坐标目标仍可能出现未定义的$\infty-\infty$，也不能只写“$\propto e^A$”掩盖该缺口。
\end{proof}

将命题\ref{prop:V5-C06-mean-field-coordinate}用于LDA之前，必须先核验当前观察词满足
$w_{mn}\in\mathcal V_{\mathrm{tr}}$且$\varphi_{k,w_{mn}}>0$对每个$k$成立。否则
$\log\varphi_{k,w_{mn}}=-\infty$，函数$A$不再满足命题中“几乎处处为实值”的假设；本章当前版本不能继续套用该命题，而应返回
$\mathtt{invalid\_input}(\mathtt{support\_failure},m,n,w_{mn})$。在上述严格正支持下，局部更新为
\begin{align}
\eta_{mnk}
&\propto \varphi_{k,w_{mn}}
\exp\!\left\{\psi(\gamma_{mk})-
\psi\!\left(\sum_{\ell=1}^{K}\gamma_{m\ell}\right)\right\},
\label{eq:V5-C06-eta-update}\\
\gamma_{mk}
&=\alpha_k+\sum_{n=1}^{N_m}\eta_{mnk}.
\label{eq:V5-C06-gamma-update}
\end{align}
每次计算$\boldsymbol\eta_{mn}$后必须对$k$归一化；在对数域使用
log-sum-exp可避免小概率下溢。式\eqref{eq:V5-C06-eta-update}和
\eqref{eq:V5-C06-gamma-update}交替至局部ELBO增量足够小。

固定全部局部变分参数，点参数变体的主题--词参数M步是归一化期望计数：
\begin{equation}
\varphi_{kv}
=\frac{\displaystyle\sum_{m=1}^{M}\sum_{n=1}^{N_m}
\eta_{mnk}\mathbf1\{w_{mn}=v\}}
{\displaystyle\sum_{v'=1}^{V}\sum_{m=1}^{M}\sum_{n=1}^{N_m}
\eta_{mnk}\mathbf1\{w_{mn}=v'\}}.
\label{eq:V5-C06-phi-update}
\end{equation}

记式\eqref{eq:V5-C06-phi-update}的分子为
$C_{kv}=\sum_{m,n}\eta_{mnk}\mathbf1\{w_{mn}=v\}$。严格正的
$\boldsymbol\varphi$与有限严格正的$\boldsymbol\gamma_m$使式
\eqref{eq:V5-C06-eta-update}归一化后有$\eta_{mnk}>0$。又因每个
$v\in\mathcal V_{\mathrm{tr}}$至少出现一次，故$C_{kv}>0$且
$C_{k\cdot}=\sum_v C_{kv}>0$，从而M步仍给出
$\varphi_{kv}=C_{kv}/C_{k\cdot}>0$。这条推导证明严格正支持是不变量；若浮点计算形成零或非有限的$C_{kv}$，那是数值支持塌缩，候选不得提交。

\begin{correctionbox}
原书对应公式给出的只是期望计数和，还不是主题--词概率。
式\eqref{eq:V5-C06-phi-update}进一步除以该主题全部词的期望计数总和。
本章选定的点参数变体不对$\boldsymbol\varphi$置先验，故式中不得临时加入
$\beta_v$或$\beta_v-1$。若要采用后验均值或MAP平滑，必须另行定义含
$\boldsymbol\varphi_k\sim\operatorname{Dir}(\boldsymbol\beta)$的目标，并同步重写ELBO、M步和评价口径。
当某主题的期望总计数为零时，式\eqref{eq:V5-C06-phi-update}无定义；当活动词表中的某个
$C_{kv}$被数值计算成零时，后续责任度和困惑度会遇到$\log0$。本章算法分别把这两类候选记为
$\texttt{numerical\_failure}(\texttt{empty\_topic},k)$与
$\texttt{numerical\_failure}(\texttt{support\_collapse},k,v)$，不进入后续Newton步，也不静默沿用旧参数或临时加平滑常数。
\end{correctionbox}

对$\boldsymbol\alpha$，令$\gamma_{m0}=\sum_k\gamma_{mk}$、
$\alpha_0=\sum_k\alpha_k$，ELBO梯度与Hessian为
\begin{align}
g_k(\boldsymbol\alpha)
&=M\{\psi(\alpha_0)-\psi(\alpha_k)\}
+\sum_{m=1}^{M}\{\psi(\gamma_{mk})-\psi(\gamma_{m0})\},
\label{eq:V5-C06-alpha-gradient}\\
H_{k\ell}(\boldsymbol\alpha)
&=M\{\psi'(\alpha_0)-\mathbf1(k=\ell)\psi'(\alpha_k)\}.
\label{eq:V5-C06-alpha-hessian}
\end{align}
Newton候选为
\begin{equation}
\boldsymbol\alpha^{\mathrm{new}}
=\boldsymbol\alpha^{\mathrm{old}}
-H(\boldsymbol\alpha^{\mathrm{old}})^{-1}
\boldsymbol g(\boldsymbol\alpha^{\mathrm{old}}).
\label{eq:V5-C06-alpha-newton}
\end{equation}
实际实现必须阻尼或线搜索，保证所有分量仍严格为正且ELBO不下降。

\begin{proposition}[受限变分EM的单调性边界]\label{prop:V5-C06-vem-monotone}
若每个E步和M步都在各自可行域内不降低同一个ELBO，则ELBO序列单调不减；
若其上有界，则ELBO数值收敛。该结论不保证参数达到全局最优，也不保证真实证据在受限平均场E步后等于ELBO。
\end{proposition}
% KNOWLEDGE-PLAN: KN-V5-C35-PROOF-007 L1
\paragraph{读前自检：证明：受限变分EM的单调性边界。}受限变分EM的单调性只在每个局部或参数子步不降低同一ELBO时成立；请比较第$t$轮前后ELBO，构造一次近似子步下降的反例，并说明此时必须拒绝候选或显式报告下降。\par

\begin{proof}
记第$t$轮参数为$(q^{(t)},\Theta^{(t)})$。坐标E步给出
\[
\mathcal L(q^{(t+1)},\Theta^{(t)})
\ge\mathcal L(q^{(t)},\Theta^{(t)}),
\]
M步给出
\[
\mathcal L(q^{(t+1)},\Theta^{(t+1)})
\ge\mathcal L(q^{(t+1)},\Theta^{(t)}).
\]
合并即得单调性；单调有界实数列收敛。由于$\mathcal Q$通常不含精确后验，
式\eqref{eq:V5-C06-elbo-identity}中的KL间隙可能严格为正；非凸性又允许不同初始化到达不同驻点。
\end{proof}

\begin{limitationbox}
有些表述把受限平均场E步写成
$\log p(\boldsymbol w\mid\Theta^{(t-1)})
=\mathcal L(q^{(t)},\Theta^{(t-1)})$。
一般只能写“$\ge$”；只有$q^{(t)}$恰为精确后验时才取等号。
因此本章把可审计结论限定为ELBO逐块单调，不夸大成每一步都精确最大化真实证据。
坐标更新达到数值稳定时，也只能报告坐标稳定点或局部驻点；多启动用于比较可行局部解，不能证明找到了变分族上的全局最优。
\end{limitationbox}


% KNOWLEDGE-PLAN: KN-V5-C35-ALGORITHM_IDEA-011 L1
\paragraph{读前自检：LDA两类推断路线总结。}LDA两类推断中，折叠Gibbs输出相关主题指派样本，变分EM输出确定性近似参数；请对照一个词元的条件更新与一个局部责任度更新，核对两者的状态、输出和诊断不同。\par
% KNOWLEDGE-PLAN: KN-V5-C35-ALGORITHM_IDEA-012 L1
\paragraph{读前自检：LDA变分EM超参数与评估。}LDA变分EM的$K,\alpha,\beta$及多启动预算必须在留出评价前冻结；请比较两个候选的留出困惑度与ELBO稳定性，核对不能用训练ELBO单独选择主题数。\par

\SLAdvancedSection{超参数、实现与评估}{sec:V5-C06-S05}
\knowledgeanchor{SLM-20-04-006}{LDA两类推断路线总结}
\knowledgeanchor{SLM-20-05-001}{LDA变分EM超参数与评估}

\phantomsection\label{struct:V5-C06-CH14}
\textbf{几何解释。}$\boldsymbol\theta_m$位于主题单纯形，表示文档是$K$个主题的凸组合；
$\boldsymbol\varphi_k$位于词汇单纯形，表示主题是$V$个词的概率组合。
模型把高维文档词频投影到低维主题单纯形，再映回词汇单纯形。若同时置换主题标签以及
$\boldsymbol\theta$、$\boldsymbol\varphi$相应坐标，观测分布不变，因此主题标签不可识别；比较两次运行需先做主题匹配。

\phantomsection\label{struct:V5-C06-CH15}
\textbf{概率解释。}给定$\boldsymbol\theta_m,\boldsymbol\varphi$时，各词位置条件独立；
边缘化共享的$\boldsymbol\theta_m$后，同文档词元相关。交换性只保证对位置作同步置换后联合概率不变，不表示词元边缘独立。
Dirichlet先验可理解为平滑伪计数：小$\alpha_0$偏好少数主题占优，大$\alpha_0$把文档主题混合拉向均值；
$\boldsymbol\beta$对主题词分布起同样作用。

\phantomsection\label{struct:V5-C06-CH16}
\textbf{优化解释。}折叠Gibbs不是直接最大化一个确定性目标，而是构造以后验为平稳分布的马尔可夫链；
变分EM则在受限分布族和模型参数之间坐标上升ELBO。前者误差来自有限链长、相关样本和混合；后者误差来自近似族、局部最优和有限迭代。二者的“收敛”不可混用。

\phantomsection\label{struct:V5-C06-CH17}
\textbf{图形辅助。}图\ref{fig:V5-C06-variational-updates}把点参数 LDA 的局部 VI 成功门与语料级 VEM 提交门分开；局部预算或失败不拥有通往 M 步的边。
% v2.7.0 figure UID: FIG-P694-01
% R2.4 compact matrix/positioning layout; no overall scaling.
\tikzset{slfig-FIG-P694-01/.style={
  font=\fontsize{9.6pt}{10.3pt}\selectfont,
  every path/.append style={line cap=round,line join=round}}}
\pgfkeysifdefined{/tikz/alt}{}{\tikzset{alt/.code={}}}
\begin{figure}[htbp]
\centering
\makebox[\linewidth][c]{%
\begin{tikzpicture}[slfig-FIG-P694-01,
  every node/.style={
    font=\fontsize{9.6pt}{10.3pt}\selectfont,
    align=center},
  box/.style={
    draw=SLDeepBlue,fill=white,rounded corners=2pt,
    line width=.78pt,inner xsep=1.6pt,inner ysep=1.6pt},
  localbox/.style={
    box,draw=SLMidBlue,fill=SLSoftTeal},
  gate/.style={
    box,double,double distance=.55pt,line width=.76pt},
  localgate/.style={
    gate,draw=SLMidBlue,fill=SLSoftTeal},
  cert/.style={
    gate,fill=SLSoftBlue},
  localcert/.style={
    localgate},
  ok/.style={
    draw=SLSelfCheckTeal,fill=SLSoftTeal,rounded corners=2pt,
    line width=.82pt,inner xsep=1.6pt,inner ysep=1.6pt},
  budget/.style={
    draw=SLAccentOrange,fill=SLWarmGray,rounded corners=2pt,
    dashed,line width=.82pt,inner xsep=1.6pt,inner ysep=1.6pt},
  fail/.style={
    draw=SLWarningRed,fill=white,rounded corners=2pt,
    dash dot,line width=.82pt,inner xsep=1.6pt,inner ysep=1.6pt},
  statusband/.style={
    draw=SLTextGray,fill=white,rounded corners=2pt,
    line width=.80pt,inner xsep=1.6pt,inner ysep=1.6pt},
  ptitle/.style={
    font=\fontsize{10.2pt}{12.2pt}\selectfont\bfseries,
    inner sep=0pt},
  main/.style={->,draw=SLDeepBlue,line width=.90pt},
  success/.style={->,draw=SLSelfCheckTeal,line width=.90pt},
  budgetedge/.style={->,draw=SLAccentOrange,line width=.82pt,dashed},
  failure/.style={->,draw=SLWarningRed,line width=.82pt,dash dot},
  >={Stealth[length=2.1mm,width=1.5mm]},
  alt={对象—关系—结论：图由内容驱动的上下两面板组成。面板a把前置及空文档门、非空初始化与同步候选、可行门及原子提交和局部双证书排成一行，下一行分别承接输入失败、数值失败、局部预算停止和唯一LocalOK；非空初始化含均匀责任度、先验加平均词数及有限初始ELBO，只有completed或converged进入LocalOK。面板b先逐文档收齐全部LocalOK，再沿第二行从右向左完成全语料计数、严格正归一、受保护Newton、完整提交和外层双证书；下一行分开外层预算候选、收敛候选和F/X非成功状态，最终仅以feasible且状态为converged或budget_stop的记录参加跨启动选择，输出保留计数、诊断和原状态。}]

% ==========================================================================
% (a) Content-driven main row.
\matrix (amain) [
  column sep=3mm,
  inner sep=0pt,
  ampersand replacement=\&
] {
  \node[localgate,text width=40mm,minimum height=29mm] (a01)
    {\textbf{A0+A1 前置／空文档门}\\
     $\mathbf w_m;\ \alpha>0$\\
     $\varphi_{k\cdot}\in\operatorname{ri}(\Delta^{V-1})$\\
     $T_{\rm loc};\quad
     \varepsilon_{\rm loc}^{\mathcal L},
     \varepsilon_{\rm loc}^{\gamma}$\\
     词表、支持、预算、容差合法？\\
     $N_m=0$：$\mathtt{completed}\to L$\\
     $N_m>0\to$ A2；非法 $\to I$};
  \&
  \node[localbox,text width=50mm,minimum height=29mm] (a2)
    {\textbf{A2 非空初始化／同步候选（$R_A$ 入口）}\\
     $\eta_{mnk}^{(0)}=1/K;\quad
     \gamma_{mk}^{(0)}=\alpha_k+N_m/K$\\
     $\mathcal L_m^{(0)}\in\mathbb R$（有限）\\
     $a_{mnk}=\log\varphi_{k,w_{mn}}+\psi(\gamma_{mk})$\\
     $\hphantom{a_{mnk}=}-\psi(\sum_\ell\gamma_{m\ell})$\\
     $\eta^+_{mnk}=e^{a_{mnk}}\big/\sum_\ell e^{a_{mn\ell}}$\\
     （LSE；同步全部 $n$）\\
     $\gamma^+_{mk}=\alpha_k+\sum_n\eta^+_{mnk}$};
  \&
  \node[localcert,text width=52mm,minimum height=29mm] (a35)
    {\textbf{A3--A5 可行门／提交／局部双证书}\\
     $F_{\rm loc}$：$\eta^+\in\operatorname{ri}(\Delta),\quad
     \gamma^+>0$\\
     $\mathcal L_m^+\in\mathbb R,\quad
     \mathcal L_m^+\ge\mathcal L_m$\\
     通过才提交；$s_{\rm done}\mathrel{+}=1,\quad
     p_{\rm done}\mathrel{+}=N_m$\\
     $\delta_{\mathcal L,m}=
     |\mathcal L_m^+-\mathcal L_m|/(1+|\mathcal L_m|)$\\
     $\delta_{\gamma,m}=
     \max_k|\gamma^+_{mk}-\gamma_{mk}|/(1+|\gamma_{mk}|)$\\
     $C_{\rm loc}$：$\delta_{\mathcal L,m}\le
     \varepsilon_{\rm loc}^{\mathcal L}$ 且\\
     $\delta_{\gamma,m}\le\varepsilon_{\rm loc}^{\gamma}$？
     否且 $s_{\rm done}<T_{\rm loc}\to R_A$};
  \\
};

\node[ptitle,text=SLMidBlue,above=1mm of amain] (atitle)
  {(a) 单文档局部 VI（$\mathrm{LL}=\mathtt{last\ legal}$）};

% (a) Four compact direct status cells; true boundary gap is 3 mm.
\matrix (astatus) [
  below=3mm of amain,
  column sep=3mm,
  inner sep=0pt,
  ampersand replacement=\&
] {
  \node[fail,text width=31mm,minimum height=12mm] (ainvalid)
    {\textbf{$I:\ \mathtt{invalid\_input}$}\\
     $\mathtt{support\_failure}$\\
     存 $\mathrm{LL}/\mathtt{diag}$；\textbf{停／禁 C/M}};
  \&
  \node[fail,text width=30mm,minimum height=12mm] (anum)
    {\textbf{$\mathtt{numerical\_failure}$}\\
     回滚 $\mathrm{LL}$；gate／位置\\
     存 diag；\textbf{停止；禁 C/M}};
  \&
  \node[budget,text width=35mm,minimum height=12mm] (alocbudget)
    {\textbf{STOP：$\mathtt{budget\_stop}$}\\
     $\neg C_{\rm loc}$；$s_{\rm done}=T_{\rm loc}$\\
     存 $\mathrm{LL}$/计数/diag；$\nexists L$；\textbf{禁 C/M}};
  \&
  \node[ok,text width=42mm,minimum height=12mm] (aok)
    {\textbf{$L:\ \mathrm{LocalOK}_m$}\\
     $(\gamma_m,\eta_m,\mathcal L_m;\ \mathtt{counts},\mathtt{diag})$\\
     $\mathtt{completed}/\mathtt{converged}$};
  \\
};

% Only adjacent main arrows; the fan-out occupies the empty 3 mm gutter.
\draw[main] (a01) -- (a2);
\draw[main] (a2) -- (a35);
\draw[failure] (a01.south west) -- (ainvalid.north);
\draw[failure] (a35.south west) -- (anum.north);
\draw[budgetedge] (a35.south) -- (alocbudget.north);
\draw[success] (a35.south east) -- (aok.north);

% ==========================================================================
% (b) Title is content-positioned below panel (a); panel backgrounds are fit
%     only after every node has its final natural bounding box.
\node[ptitle,text=SLDeepBlue,below=3mm of astatus] (btitle)
  {(b) 语料级外层 VEM};

\matrix (btop) [
  below=1mm of btitle,
  column sep=3mm,
  inner sep=0pt,
  ampersand replacement=\&
] {
  \node[box,text width=52mm,minimum height=19mm] (b0)
    {\textbf{B0 全局门／多启动}\\
     $\mathcal D,K;\ \alpha,\varphi>0;\quad T_{\rm loc},T_{\rm out}$\\
     $\varepsilon_{\rm out}^{\mathcal L,\Theta};\quad
     s=1{:}R_{\rm start};\ J_{\rm ls}$\\
     随机源失败：$\mathtt{random\_source\_failure}$\\
     其余前置门否：$\mathtt{invalid\_input}$\\
     通过 $\to r_{\rm done}=0$};
  \&
  \node[box,text width=54mm,minimum height=19mm] (b1)
    {\textbf{B1 全部文档（$R$ 入口）}\\
     对 $m=1,\ldots,M$ 调用面板 (a)\\
     收集 $(\mathrm{LocalOK}_m,\mathtt{status}_m,
     \mathcal L_m,\mathtt{diag}_m)$\\
     局部 $\mathtt{budget\_stop}\to X$\\
     局部 failure $\to F$};
  \&
  \node[gate,text width=28mm,minimum height=19mm] (b2)
    {\textbf{B2 全收齐门}\\
     $\forall m:\ \mathrm{LocalOK}_m$ 到齐？\\
     是 $\to$ B3\\
     否 $\to F/X$（按原 status）};
  \\
};

% Right-to-left second row: B3 -> B4 -> B5/B6.
\matrix (bmain) [
  below=3mm of btop,
  column sep=3mm,
  inner sep=0pt,
  ampersand replacement=\&
] {
  \node[cert,text width=56mm,minimum height=29mm] (b56)
     {\textbf{B5--B6 提交／证书（$R:\ r<T_{\rm out}$）}\\
     $\Theta^+=(\alpha^+,\varphi^+,
     \{\gamma_m^+,\eta_m^+\}_{m=1}^{M},\mathcal L^+)$\\
     可行、有限且不下降才提交；
     $r_{\rm done}\mathrel{+}=1$\\
     $\delta_{\mathcal L}=
     |\mathcal L^+-\mathcal L|/(1+|\mathcal L|)$\\
     $\delta_\alpha=
     \max_k|\alpha_k^+-\alpha_k|/(1+|\alpha_k|)$\\
     $\delta_\varphi=\max_{k,v}
     |\varphi^+_{kv}-\varphi_{kv}|$\\
     $\delta_\Theta=\max\{\delta_\alpha,\delta_\varphi\}$\\
     $C_{\rm out}$：
     $\delta_{\mathcal L}\le\varepsilon_{\rm out}^{\mathcal L}
     \land\delta_\Theta\le\varepsilon_{\rm out}^{\Theta}$？};
  \&
  \node[cert,text width=40mm,minimum height=29mm,inner ysep=2.4pt] (b4)
    {\textbf{B4 受保护 Newton}\\
     \textbf{（$N$ 入口）}\\
     $H(\alpha)\Delta=g(\alpha)$\\
     $\alpha^+=\alpha-\rho\Delta$\\
     $\rho\in\{1,\tfrac12,\tfrac14,\ldots\}$；
     $j\le J_{\rm ls}$\\
     $G_N$：$\min_k\alpha_k^+>0$\\
     且新 ELBO 有限、不下降\\
     拒绝且 $j<J_{\rm ls}$：
     $\rho\leftarrow\rho/2\to N$\\
     $j=J_{\rm ls}$：
     $\mathtt{line\_search\_failed}\to F$};
  \&
  \node[box,text width=42mm,minimum height=29mm] (b3)
    {\textbf{B3 全语料计数 $\to\varphi^+$}\\
     $C^+_{kv}=\sum_{m=1}^{M}\sum_{n=1}^{N_m}$\\
     $\eta_{mnk}\mathbf1\{w_{mn}=v\}$\\
     $G_M$：计数有限；$C^+_{k\cdot}>0$\\
     且每个 $C^+_{kv}>0$\\
     $\varphi^+_{kv}=C^+_{kv}/C^+_{k\cdot}>0$\\
     $\sum_v\varphi^+_{kv}=1$\\
     失败 $\to F:\mathtt{numerical\_failure}$};
  \\
};

% Compact two-row terminal table.  Its top row is addressable by all exits;
% the full-width selection row is positioned from the completed top bbox.
\matrix (bterminal) [
  below=3mm of bmain,
  column sep=3mm,
  inner sep=0pt,
  ampersand replacement=\&
] {
  \node[budget,text width=48mm,minimum height=12mm] (bbudget)
    {\textbf{STOP：$\mathtt{budget\_stop}$（未收敛候选）}\\
     $\neg C_{\rm out},\ r=T_{\rm out}$（本启动）\\
     可行；存 $\mathrm{LL}$/计数/diag（原 status）};
  \&
  \node[ok,text width=36mm,minimum height=12mm] (bconv)
    {$\checkmark$ \textbf{$\mathtt{converged}$ 候选}\\
     $C_{\rm out}$；$\mathtt{feasible}=\mathtt{true}$\\
     存 $\mathtt{record}[s]$（原 status）};
  \&
  \node[statusband,text width=58mm,minimum height=12mm] (bfail)
    {\textbf{$F/X$：非成功；禁 C/M；$\notin\mathcal S_{\rm acc}$}\\
     $F$：保留入口的原失败 status\\
     $X$：局部 $\mathtt{budget\_stop}$（原 status）\\
     $\mathtt{feasible}=\mathtt{false}$；仅存 status／$\mathrm{LL}$/diag};
  \\
};

\node[box,text width=148mm,minimum height=15mm,fill=SLSoftTeal,
  font=\fontsize{9.6pt}{9.9pt}\selectfont,
  below=2mm of bterminal] (bselect)
  {\textbf{跨启动选择／输出}\quad $q_s:=\mathtt{record}[s]$\\
   $\mathcal S_{\rm acc}=\{s:q_s.\mathtt{feasible}=\mathtt{true},\quad
   q_s.\mathtt{status}\in\{\mathtt{converged},\mathtt{budget\_stop}\}\}$\\
   空集：仅返回全部 records；非空：按最终 ELBO／冻结并列规则返回
   $(\widehat\alpha,\widehat\varphi,
   \{\widehat\gamma_m,\widehat\eta_m\},\mathcal L)$\\
   $s_{\rm done},r_{\rm done},m_{\rm done},k_{\rm done},j_{\rm done}$；
   diagnostics；原 status：
   \textcolor{SLSelfCheckTeal}{$\checkmark\mathtt{converged}$}／
   \textcolor{SLAccentOrange}{STOP $\mathtt{budget\_stop}$（未收敛）}};

% Main edges are only between adjacent nodes. Cross-row edges stay in the
% 3 mm boundary gutters and preserve left-to-right endpoint ordering.
\draw[main] (b0) -- (b1);
\draw[main] (b1) -- (b2);
\draw[main] (b2.south) -- (b3.north);
\draw[main] (b3) -- (b4);
\draw[main] (b4) -- (b56);

% B3 and B4 each have a real, short failure edge to the right-hand F/X cell.
\draw[failure] (b3.south) -- (bfail.north east);
\draw[failure] (b4.south east) -- (bfail.north west);
\draw[success] (b56.south east) -- (bconv.north);
\draw[budgetedge] (b56.south west) -- (bbudget.north);

\draw[budgetedge] (bbudget.south) -- (bselect.north west);
\draw[success] (bconv.south) -- (bselect.north);

% Content-derived backgrounds are created last, after every real node bbox.
\begin{scope}[on background layer]
  \node[
    draw=SLMidBlue,fill=SLSoftTeal,fill opacity=.24,
    line width=.80pt,rounded corners=4pt,
    fit=(atitle)(amain)(astatus),
    inner xsep=2mm,inner ysep=1mm] (apanel) {};
  \node[
    draw=SLDeepBlue,fill=SLSoftBlue,fill opacity=.22,
    line width=.80pt,rounded corners=4pt,
    fit=(btitle)(btop)(bmain)(bterminal)(bselect),
    inner xsep=2mm,inner ysep=1mm] (bpanel) {};
\end{scope}

\end{tikzpicture}%
}
\captionsetup{width=.97\linewidth}
\caption{点参数 LDA 两层变分 EM：仅 $\mathtt{completed}/\mathtt{converged}$
形成 $\mathrm{LocalOK}_m$，全部 $m$ 到齐后才执行
$C\to\varphi\to$ 受保护 Newton；可行的 $\mathtt{budget\_stop}$
保留原 status，作为未收敛候选进入 $\mathcal S_{\rm acc}$，
其余非成功状态止于记录。}
\label{fig:V5-C06-variational-updates}
\end{figure}

\FloatBarrier
\textbf{读图检查。}先沿 (a) 逐文档确认只有 $\mathtt{completed}/\mathtt{converged}$ 形成 $\mathrm{LocalOK}_m$，局部预算或失败不进 $C$/M 步；再沿 (b) 核对全部 $\mathrm{LocalOK}_m$ 到齐后才可进入 $C$、$\varphi$ 与受保护 Newton。外层 $\mathtt{budget\_stop}$ 只有在该启动仍 $\mathtt{feasible}$ 时才以“预算候选／未收敛”进入 $\mathcal S_{\rm acc}$，其余失败终点只留记录。

图\ref{fig:V5-C06-method-comparison}把完整Bayes折叠Gibbs与无
$\boldsymbol\beta$先验的点参数变分EM按模型目标、局部/全局变量、核心更新、
输出与适用性、风险与诊断逐项并列。
% v2.7.0 figure UID: FIG-P695-01
% Two-model comparison: fixed five-row grid, explicit model boundary, no connectors.
\tikzset{slfig-FIG-P695-01/.style={font=\fontsize{9.6pt}{11.6pt}\selectfont,every path/.append style={line cap=round,line join=round}}}
\pgfkeysifdefined{/tikz/alt}{}{\tikzset{alt/.code={}}}
\begin{figure}[htbp]
\centering
\begin{tikzpicture}[slfig-FIG-P695-01,
  every node/.style={font=\fontsize{9.6pt}{11.6pt}\selectfont,align=left,
    inner xsep=2.0mm,inner ysep=1.3mm,outer sep=0pt},
  head/.style={draw=SLDeepBlue,fill=SLDeepBlue,text=white,
    rounded corners=1.5pt,minimum height=10mm,align=center,
    font=\fontsize{10.2pt}{12.2pt}\selectfont\bfseries,line width=0.84pt},
  key/.style={draw=SLRuleGray,fill=SLSoftGray,
    rounded corners=1.2pt,text width=25mm,align=center,
    font=\fontsize{9.6pt}{11.6pt}\selectfont\bfseries,line width=0.74pt},
  gibbs/.style={draw=SLMidBlue,fill=SLSoftBlue,
    rounded corners=1.2pt,text width=55.5mm,align=left,
    line width=0.78pt},
  variational/.style={draw=SLRuleGray,dashed,fill=white,
    rounded corners=1.2pt,text width=55.5mm,align=left,
    line width=0.78pt},
  modelrow/.style={minimum height=18mm},
  variablerow/.style={minimum height=16mm},
  updaterow/.style={minimum height=23mm},
  outputrow/.style={minimum height=18mm},
  riskrow/.style={minimum height=23mm},
  alt={对象—关系—结论：左列是完整Bayes LDA的折叠Gibbs，积分消去随机的文档主题向量和主题词向量，以主题指派及计数为链状态并输出相关后验样本；右列是无主题词Dirichlet先验的点参数LDA变分EM，局部E步只固定当前主题词参数并更新每篇文档的gamma与eta，全局再以期望计数更新varphi、以受保护Newton更新alpha。前者适合需要后验不确定性或多峰信息的情形，后者适合大语料、批量化或预算敏感情形。同一点参数模型的可行块更新可使ELBO逐块非降，但ELBO不等于真实证据；公平比较须冻结相同切分、词表、预算与评价口径，并同时登记模型差异和推断差异。}]
\node[head,text width=25mm] at (-6.35,6.60) {比较维度};
\node[head,text width=55.5mm] at (-1.55,6.60)
  {折叠 Gibbs\\（完整 Bayes LDA）};
\node[head,text width=55.5mm] at (4.83,6.60)
  {平均场变分 EM\\（点参数 LDA）};

\node[key,modelrow] at (-6.35,4.80) {模型目标};
\node[gibbs,modelrow] at (-1.55,4.80)
  {\textbf{完整 Bayes LDA：}\\
   $\boldsymbol\theta_m\sim\operatorname{Dir}(\boldsymbol\alpha)$，
   $\boldsymbol\varphi_k\sim\operatorname{Dir}(\boldsymbol\beta)$；\\
   折叠二者，抽样$p(\boldsymbol z\mid\boldsymbol w,\boldsymbol\alpha,\boldsymbol\beta)$。};
\node[variational,modelrow] at (4.83,4.80)
  {\textbf{点参数 LDA：无 $\boldsymbol\beta$ 先验。}\\
   $\boldsymbol\varphi_k$为待估参数；最大化\\
   同一点参数模型的ELBO。};

\node[key,variablerow] at (-6.35,2.85) {局部/全局\\变量};
\node[gibbs,variablerow] at (-1.55,2.85)
  {局部链状态：$z_i$；共享计数：\\
   $n_{mk},n_{kv},n_{k\cdot}$；$\boldsymbol\theta,\boldsymbol\varphi$已折叠。};
\node[variational,variablerow] at (4.83,2.85)
  {局部E步固定\textbf{当前}$\boldsymbol\varphi$：\\
   $(\boldsymbol\gamma_m,\boldsymbol\eta_m)$；全局：
   $(\boldsymbol\varphi,\boldsymbol\alpha)$。};

\node[key,updaterow] at (-6.35,0.65) {核心更新};
\node[gibbs,updaterow] at (-1.55,0.65)
  {\textbf{留一计数满条件：}\\[-0.3mm]
   $\displaystyle p(z_i=k\mid\cdot)\propto
   \frac{n_{kv}^{-i}+\beta_v}{n_{k\cdot}^{-i}+\beta_0}
   (n_{mk}^{-i}+\alpha_k)$；\\
   执行“减--算--抽--加”。};
\node[variational,updaterow] at (4.83,0.65)
  {局部$\boldsymbol\eta_m\leftrightarrow\boldsymbol\gamma_m$；
   全部文档成功后\\
   $C_{kv}=\sum_{m,n}\eta_{mnk}\mathbf1\{w_{mn}=v\}$；\\
   $\varphi_{kv}=C_{kv}/C_{k\cdot}$；
   $\boldsymbol\alpha$用受保护Newton，\\
   候选须在正域且ELBO不下降。};

\node[key,outputrow] at (-6.35,-1.65) {输出、优势\\与适用性};
\node[gibbs,outputrow] at (-1.55,-1.65)
  {输出相关后验样本；保留后验不确定性\\
   与多峰信息。稀疏计数实现自然；\\
   适于需要不确定性或多峰信息时。};
\node[variational,outputrow] at (4.83,-1.65)
  {输出$(\boldsymbol\varphi,\boldsymbol\alpha,\{q_m\})$点参数/近似；\\
   确定且易批量化；适于大语料、\\
   批量处理或预算敏感情形。};

\node[key,riskrow] at (-6.35,-3.95) {风险与诊断};
\node[gibbs,riskrow] at (-1.55,-3.95)
  {烧入不等于收敛；检查混合、自相关/ESS、\\
   多链主题匹配后的一致性；\\
   单链样本轨迹不能充当收敛证书。};
\node[variational,riskrow] at (4.83,-3.95)
  {近似族与局部最优；检查多启动、局部未收敛、\\
   空主题、线搜索与数值失败。同一点参数模型的\\
   可行块更新使ELBO逐块非降；\\
   ELBO不等于真实证据。};

\node[draw=SLDeepBlue,fill=SLSoftBlue,rounded corners=2pt,
  text width=151mm,minimum height=11mm,align=center,line width=0.84pt]
  at (0,-5.90)
  {\textbf{公平比较：}冻结同一训练/验证/测试切分、词表与预处理、$K$、计算预算和评价口径；
   同时登记模型差异与推断差异。};
\end{tikzpicture}
\captionsetup{width=.94\linewidth}
\caption{完整Bayes折叠Gibbs与无$\boldsymbol\beta$先验的点参数变分EM：二者模型目标、状态、更新、输出和诊断均不同，公平比较须同时登记模型与推断差异。}
\label{fig:V5-C06-method-comparison}
\end{figure}

\FloatBarrier
\noindent\textbf{读图检查。}先确认左列积分消去
$\boldsymbol\theta,\boldsymbol\varphi$并以$\boldsymbol z$和计数为链状态，
右列只在局部E步固定当前$\boldsymbol\varphi$，再由全局期望计数与受保护Newton分别更新
$\boldsymbol\varphi,\boldsymbol\alpha$；最后只在相同训练/验证/测试切分、
词表与预处理、主题数、计算预算和评价口径下比较，并同时登记模型差异与推断差异。

\phantomsection\label{struct:V5-C06-CH18}
\textbf{小型数值例子。}设$K=2$，当前词$v$的留一计数为
\[
(n_{1v}^{-i},n_{2v}^{-i})=(3,1),\quad
(n_{1\cdot}^{-i},n_{2\cdot}^{-i})=(8,6),\quad
(n_{m1}^{-i},n_{m2}^{-i})=(2,4),
\]
并取对称$\alpha=(1,1)$、每词$\beta_v=1$、$V=5$。
式\eqref{eq:V5-C06-gibbs-conditional}给出未归一化权重
\[
r_1=\frac4{13}\frac3{8}=\frac3{26},\qquad
r_2=\frac2{11}\frac5{8}=\frac5{44}.
\]
归一化后$p(z_i=1)=66/131$、$p(z_i=2)=65/131$。两个局部因子接近，说明不能只看“当前词更像哪个主题”，还要同时看文档主题计数。

对一次变分坐标更新，若
$\boldsymbol\gamma=(3,2)$且该词在两主题下概率为$(0.6,0.2)$，则
\[
\psi(3)-\psi(5)=-\frac7{12},\qquad
\psi(2)-\psi(5)=-\frac{13}{12}.
\]
未归一化责任度约为$(0.3348,0.0677)$，归一化为$(0.832,0.168)$。


\phantomsection\label{struct:V5-C06-CH19}
\textbf{伪代码。}主线只保留两条推断算法：完整Bayes折叠Gibbs与点参数LDA变分EM；算法\ref{alg:V5-C06-generate}仅用于说明生成方向。通用变分EM和单文档局部更新改作稳健实现细节条带，避免把模型路线、通用模板和子程序误读为五个并列选择。

\phantomsection\label{struct:V5-C06-CH20}
\textbf{输入与输出。}生成算法输入维数、长度、先验和随机数，输出完整模拟语料；
Gibbs输入语料、$K,\boldsymbol\alpha,\boldsymbol\beta$及扫描计划，输出保留样本、后验均值和诊断；
通用变分EM输入模型、数据、可行族和两个块优化器，输出$q$与模型参数；
局部变分推断输入一篇文档、固定全局参数和容差，输出$\boldsymbol\gamma_m,\boldsymbol\eta_m$及局部ELBO；
完整变分EM输入语料和初始化，输出全局参数、全部局部参数和ELBO轨迹。

\phantomsection\label{struct:V5-C06-CH21}
\textbf{参数含义。}$K$控制表示维度而非已知真主题数；$\boldsymbol\alpha$控制文档主题稀疏与均值；
$\boldsymbol\beta$只属于完整Bayes/Gibbs模型并控制主题词先验，点参数变分EM不使用它；
$T_{\mathrm B}$是燃烧扫描数，$T_{\mathrm S}$是留样扫描数，$L$是抽样间隔；
$\varepsilon_{\mathrm{loc}},\varepsilon_{\mathrm{out}}$分别是局部和外层ELBO相对增量阈值；
$T_{\mathrm{loc}},T_{\mathrm{out}}$是硬迭代上限；$R_{\mathrm{chain}}$是独立Gibbs链数，
$R_{\mathrm{start}}$是变分随机启动数，$J_{\mathrm{ls}}$是每次Newton更新的最大回溯候选数。

\phantomsection\label{struct:V5-C06-CH22}
\begin{implementationbox}
Gibbs可按固定种子均匀指派主题并一次性建立两张计数表；必须核验表的行列和与$\boldsymbol z$一致。
变分EM可用多次随机单纯形初始化、确定性分块初始化或外部可复现实验初始化；每篇文档取
$\eta_{mnk}=1/K$、$\gamma_{mk}=\alpha_k+N_m/K$是合法起点。空文档不更新主题责任度，只保留先验$\boldsymbol\gamma_m=\boldsymbol\alpha$。

\phantomsection\label{struct:V5-C06-CH23}
\textbf{循环变量与更新顺序。}Gibbs按扫描$t$、文档$m$、位置$n$、候选主题$k$循环，严格执行“减旧计数--算$K$个权重--归一化--抽新主题--加新计数”。
局部变分按迭代$s$，先更新所有$\boldsymbol\eta_{mn}$，再汇总更新$\boldsymbol\gamma_m$。
外层变分EM按轮次$r$，先对所有文档完成局部E步，再汇总期望计数更新$\boldsymbol\varphi$，最后以受保护Newton步更新$\boldsymbol\alpha$。

\phantomsection\label{struct:V5-C06-CH24}
\textbf{判断、停止、返回与异常。}任何非法词索引、非正超参数、非有限概率、计数不守恒或归一化失败都使过程立即停止并指出位置。
Gibbs在完成预设燃烧与留样计划后停止；若达到预算仍未满足诊断条件，应说明当前样本尚不足以支持稳定推断。变分循环同时检查ELBO不下降、相对增量和硬上限；若Newton候选越过正域则回溯。
结果应同时给出实际迭代数、目标轨迹和初始化方式。
\end{implementationbox}

\phantomsection\label{struct:V5-C06-CH25}
\Needspace{3\baselineskip}
\textbf{正确性与收敛条件。}Gibbs更新由命题\ref{prop:V5-C06-gibbs-conditional}保证；在有限状态、不可约、非周期等条件下链以目标后验为平稳分布，但有限样本仍有Monte Carlo误差。
变分更新由命题\ref{prop:V5-C06-mean-field-coordinate}与
\ref{prop:V5-C06-vem-monotone}保证同一个点参数模型ELBO逐块非降。
零概率锁死、主题空置、坏初始化、过严/过松容差、数值下溢或Newton越界都可能导致数值失败或预算停止；这些状态都不得写成数学收敛。

\phantomsection\label{struct:V5-C06-CH26}
\textbf{训练时间复杂度与预测时间复杂度。}令$I=\sum_mN_m$。稠密实现的一次Gibbs扫描计算$K$个候选，时间为$O(IK)$；
$R_{\mathrm{chain}}$条独立链的总训练时间为
$O\!\left(R_{\mathrm{chain}}(T_{\mathrm B}+LT_{\mathrm S})IK\right)$。
一次局部变分全语料扫描为$O(IK)$，一次M步汇总亦为$O(IK+KV)$。
若每轮Newton最多试$J_{\mathrm{ls}}$个回溯候选，并利用式\eqref{eq:V5-C06-alpha-hessian}
“对角矩阵加秩一矩阵”的结构在线性时间内求解，则$R_{\mathrm{start}}$次启动的总时间为
\[
O\!\left(R_{\mathrm{start}}T_{\mathrm{out}}
\left[T_{\mathrm{loc}}IK+IK+KV+J_{\mathrm{ls}}K\right]\right).
\]
若每个Newton候选改用普通稠密线性求解，则最后一项改为
$J_{\mathrm{ls}}K^3$。因此不能把多启动或回溯候选数藏进常数。
新文档折入固定主题模型需$O(T_{\mathrm{loc}}N_mK)$，逐词预测另需$O(N_mK)$。

\phantomsection\label{struct:V5-C06-CH27}
\textbf{空间复杂度与复杂度假设。}Gibbs当前链的计数与指派工作内存为
$O(KV+MK+I)$；若把$T_{\mathrm S}$个完整$\boldsymbol z$样本全部驻留内存，还须加
$O(T_{\mathrm S}I)$，而逐样本流式落盘或在线汇总可保持前一工作内存界。
批量变分若把全部责任度作为驻留输出需$O(IK+KV+MK)$；只有在
$\boldsymbol\eta$逐文档流式落盘或不保留时，E步工作内存才可降到
$O(KV+MK+K\max_mN_m)$。上述界把固定精度加法、digamma和一次类别抽样视为$O(1)$；稀疏词表、别名采样和并行同步会改变常数或参数化界，不能无条件宣称更低复杂度。

\begingroup\SetAlgoLined
% KNOWLEDGE-PLAN: KN-V5-C35-ALGORITHM_IDEA-013 L2

% KNOWLEDGE-PLAN: KN-V5-C35-ALGORITHM_IDEA-014 L2
\begin{keypointbox}[title={LDA折叠Gibbs抽样：五步阅读}]
\textbf{输入。}写出数据对象、固定参数与必须成立的条件。\par
\textbf{初始化。}标明主状态、计数器和第一个合法状态。\par
\textbf{核心更新。}只手算一次从旧状态到候选状态的更新，并指出何时提交。\par
\textbf{数学停止。}写出能直接核验的停止证书；预算耗尽不等同于收敛。\par
\textbf{输出。}列出返回对象、实际计数、中心状态和用于复核的诊断。
\end{keypointbox}

\begin{AlgorithmContract}{
  input={合法语料、主题数、严格正先验、完整初始指派、固定词元顺序、扫描预算$T$与随机源},
  output={最后合法指派/计数、完整扫描链、$t_{\rm done}$、成功词元步$i_{\rm done}$、随机调用$q_{\rm done}$、状态与最终诊断},
  preconditions={$K\ge1,T\in\mathbb N_0$；词索引与初始主题索引合法；计数可由指派一致重建；随机类别接口合法},
  initialization={由$z^{(0)}$重建并核对$n_{mk},n_{kv},n_{k\cdot}$，置$t_{\rm done}=i_{\rm done}=q_{\rm done}=0$},
  loop={按固定扫描和词元次序，对每词元先临时减旧计数、形成满条件、抽新主题并恢复新计数},
  update={减计数、抽样与加计数组成原子单词元事务；任一步失败均回滚旧指派/计数；整扫描完成后才追加链},
  domain={全部计数为非负整数且边际一致；条件概率有限非负、和为正并归一；指派索引合法},
  failure={语料/初值/预算非法返回$\mathtt{invalid\_input}$；类别随机源失败返回$\mathtt{random\_source\_failure}$；计数、条件归一化或索引异常返回$\mathtt{numerical\_failure}$并原子回滚},
  stop={完成$T$个完整扫描、$T=0$或首次词元失败时停止；达到扫描上限只表示$\mathtt{completed}$的固定预算执行},
  budget={至多$TI$个词元条件更新和随机调用，$I$为语料词元数},
  status={$\mathtt{completed}$、$\mathtt{invalid\_input}$、$\mathtt{numerical\_failure}$、$\mathtt{random\_source\_failure}$},
  iterations={$t_{\rm done}$计完整提交扫描，$i_{\rm done}$计成功词元事务，$q_{\rm done}$计实际类别随机调用},
  diagnostics={失败$(t,i,\mathrm{stage})$、回滚确认、最小条件归一化常数、计数不变量和计划/实际扫描},
  complexity={$O(TIK)$时间；当前链工作空间$O(KV+MK+I)$，保存全部扫描另需$O(TI)$}
}
% v2.6 layout: former forced clear removed; natural pagination.
\begin{algorithm}[tbp]
\caption{LDA折叠Gibbs抽样}\label{alg:V5-C06-gibbs}
\KwIn{语料$\boldsymbol w$；主题数$K$；严格正$\boldsymbol\alpha,\boldsymbol\beta$；初始指派$\boldsymbol z^{(0)}$；扫描数$T$}
\KwOut{最后合法指派/计数与完整扫描链；$t_{\rm done},i_{\rm done},q_{\rm done}$；$\mathtt{status}$与最终诊断}
若语料、$K,T$、先验、初始指派或随机接口非法，则返回空结果、计数$0$、$\mathtt{invalid\_input}$及字段诊断\;
由$z^{(0)}$建立并核对计数，置$t_{\rm done}=i_{\rm done}=q_{\rm done}=0$；重建不变量失败则返回初始指派、$\mathtt{numerical\_failure}$及初始化诊断\;
若$T=0$，则返回初始指派/计数、空扫描链、$\mathtt{completed}$及零预算诊断\;
\For{扫描$t\leftarrow1$ \KwTo $T$}{
  \For{$i$依预定词元次序遍历，且$i=(m,n),\ v=w_i$}{
    保存旧指派和三项计数，再在临时事务中减一；若将产生负计数则回滚并返回$\mathtt{numerical\_failure}$及$(t,i)$计数诊断\;
    计算并归一化$K$个条件概率；非有限、负值或总和非正时回滚并返回最后合法状态、$\mathtt{numerical\_failure}$及条件诊断\;
    调用随机源并增加$q_{\rm done}$；调用失败或主题索引非法时回滚，分别返回$\mathtt{random\_source\_failure}$或$\mathtt{numerical\_failure}$及抽样诊断\;
    原子提交新指派及三项加一并增加$i_{\rm done}$\;
  }
  原子追加$z^{(t)}\leftarrow z$并令$t_{\rm done}\leftarrow t$\;
}
返回最后合法指派/计数、完整链、实际计数、$\mathtt{completed}$及计数不变量最终诊断\;
\end{algorithm}
\end{AlgorithmContract}

% KNOWLEDGE-PLAN: KN-V5-C35-ALGORITHM_IDEA-015 L2
\begin{keypointbox}[title={算法执行契约：五步阅读}]
\textbf{输入。}写出数据对象、固定参数与必须成立的条件。\par
\textbf{初始化。}标明主状态、计数器和第一个合法状态。\par
\textbf{核心更新。}只手算一次从旧状态到候选状态的更新，并指出何时提交。\par
\textbf{数学停止。}写出能直接核验的停止证书；预算耗尽不等同于收敛。\par
\textbf{输出。}列出返回对象、实际计数、中心状态和用于复核的诊断。
\end{keypointbox}

\begin{AlgorithmContract}{
  input={数据、联合模型、变分族、E/M更新接口、可行初值、外层预算$T_{\max}$、ELBO/残差容差},
  output={最后合法$q,\Theta,\mathcal L$、ELBO轨迹、提交轮数$t_{\rm done}$、E步数$e_{\rm done}$、M步数$m_{\rm done}$、状态与最终诊断},
  preconditions={$T_{\max}\in\mathbb N_0$；初值可行；ELBO和一阶残差可稳定计算；E/M接口承诺同一模型目标},
  initialization={$q=q^{(0)},\Theta=\Theta^{(0)}$，新鲜计算有限ELBO/残差，$t_{\rm done}=e_{\rm done}=m_{\rm done}=0$},
  loop={每轮在临时对象中先固定旧$\Theta$完成E步，再固定新$q^+$完成M步并复算同一ELBO},
  update={$q^+,\Theta^+,\mathcal L^+$、可行性和单调门全部通过后才原子提交完整轮},
  domain={$q\in\mathcal Q$、$\Theta$在参数域；ELBO/残差有限；提交轨迹不下降超过舍入容差},
  failure={接口/初值/预算非法返回$\mathtt{invalid\_input}$；E/M更新、ELBO、可行性或单调门失败返回$\mathtt{numerical\_failure}$并回滚到上一合法轮},
  stop={ELBO相对增量与一阶残差双证书、$T_{\max}$轮耗尽或首次失败时停止},
  budget={至多$T_{\max}$个完整E--M提交轮，各轮各一次登记E步和M步接口调用},
  status={$\mathtt{converged}$、$\mathtt{budget\_stop}$、$\mathtt{invalid\_input}$、$\mathtt{numerical\_failure}$},
  iterations={$t_{\rm done}$计完整提交轮，$e_{\rm done}$和$m_{\rm done}$计实际成功完成的E/M接口调用},
  diagnostics={初末ELBO、相对增量、一阶残差、失败$(t,\mathrm{E/M/ELBO})$、回滚确认及预算余量},
  complexity={总时间为实际E/M调用成本之和；至少保存最后合法与当前候选两套状态}
}
\phantomsection\label{alg:V5-C06-generic-vem}
\SLRobustImplementationStrip{通用变分EM模板：每轮在临时对象上依次完成E块与M块，复算同一ELBO及一阶残差；只有可行性、单调门和有限性全通过才原子提交。双证书满足报$\mathtt{converged}$，预算耗尽报$\mathtt{budget\_stop}$，接口或单调门失败回滚并报$\mathtt{numerical\_failure}$。}
\end{AlgorithmContract}

% KNOWLEDGE-PLAN: KN-V5-C35-ALGORITHM_IDEA-016 L2
\begin{keypointbox}[title={算法执行契约：五步阅读}]
\textbf{输入。}写出数据对象、固定参数与必须成立的条件。\par
\textbf{初始化。}标明主状态、计数器和第一个合法状态。\par
\textbf{核心更新。}只手算一次从旧状态到候选状态的更新，并指出何时提交。\par
\textbf{数学停止。}写出能直接核验的停止证书；预算耗尽不等同于收敛。\par
\textbf{输出。}列出返回对象、实际计数、中心状态和用于复核的诊断。
\end{keypointbox}

\begin{AlgorithmContract}{
  input={单文档词索引、训练活动词表、固定严格正$\alpha$与在活动词表上严格正的主题--词概率$\varphi$、局部预算$T_{\rm loc}$及ELBO/参数容差},
  output={最后合法$\gamma_m,\eta_m$与局部ELBO、提交轮$s_{\rm done}$、责任度更新$p_{\rm done}$、状态与最终诊断},
  preconditions={词索引属于训练活动词表，$\alpha$有限严格正，$\varphi$逐主题归一且每个目标词在每个主题下严格为正，预算非负、容差合法},
  initialization={非空文档用均匀$\eta$及$\gamma=\alpha+N_m/K$并计算基线ELBO；$s_{\rm done}=p_{\rm done}=0$},
  loop={每轮在临时表中逐词以log-sum-exp更新责任度，再同步形成$\gamma^+$与局部ELBO},
  update={全篇$\eta^+,\gamma^+$、ELBO和单纯形检查全部通过后才原子提交完整局部轮},
  domain={$\eta_{mn}\in\operatorname{ri}(\Delta^{K-1})$、$\gamma_{mk}>0$、目标词混合概率严格为正且ELBO有限},
  failure={词表外项或任一目标词概率非正返回$\mathtt{invalid\_input}(\mathtt{support\_failure})$；其他接口/词索引/概率模型非法返回$\mathtt{invalid\_input}$；责任度归一化、digamma、ELBO或单调门失败返回$\mathtt{numerical\_failure}$并回滚},
  stop={空文档直接$\mathtt{completed}$；非空文档双容差证书、$T_{\rm loc}$轮耗尽或首次失败时停止},
  budget={非空文档至多$T_{\rm loc}$个完整局部轮和$T_{\rm loc}N_m$次责任度向量更新},
  status={$\mathtt{completed}$、$\mathtt{converged}$、$\mathtt{budget\_stop}$、$\mathtt{invalid\_input}$、$\mathtt{numerical\_failure}$},
  iterations={$s_{\rm done}$计完整提交局部轮，$p_{\rm done}$计成功形成的词责任度候选},
  diagnostics={空文档分支、初末ELBO、相对增量、最大参数变化、失败$(s,n)$与回滚、预算余量},
  complexity={$O(s_{\rm done}N_mK)$时间、$O(N_mK+K)$空间}
}
\phantomsection\label{alg:V5-C06-local-vi}
\SLRobustImplementationStrip{单文档局部子程序：空文档返回$\gamma_m=\alpha$；非空文档在对数域同步更新式\eqref{eq:V5-C06-eta-update}与\eqref{eq:V5-C06-gamma-update}，整篇候选通过ELBO与参数双门后才提交。词表支持失败报$\mathtt{invalid\_input}$，预算耗尽报$\mathtt{budget\_stop}$，数值或单调门失败回滚。}
\end{AlgorithmContract}

% KNOWLEDGE-PLAN: KN-V5-C35-ALGORITHM_IDEA-017 L1
\paragraph{读前自检：点参数LDA变分EM（多启动）：执行契约。}点参数LDA变分EM（多启动）输入“训练语料及活动词表、主题数$K$、严格正初始$\alpha,\varphi$、局部/外层双预算和容差、可复现启动表及Newton回溯预算”，条件“语料词索引和维数合法，活动词表每项至少出现一次，$K,R_{\rm start}\ge1$，初值在整个活动词表上严格正归一，全部预算为非负整数且容差/并列规则冻结”；请执行“完整局部推断、临时主题矩阵、临时$\alpha$和新ELBO全部通过归一化/单调门禁后，才原子提交一次外层轮”，以“每启动满足外层ELBO/参数双证书、外层预算耗尽或首次启动内失败”判停。\par
% KNOWLEDGE-PLAN: KN-V5-C35-ALGORITHM_IDEA-018 L2
\begin{keypointbox}[title={点参数LDA变分EM（多启动）：五步阅读}]
\textbf{输入。}写出数据对象、固定参数与必须成立的条件。\par
\textbf{初始化。}标明主状态、计数器和第一个合法状态。\par
\textbf{核心更新。}只手算一次从旧状态到候选状态的更新，并指出何时提交。\par
\textbf{数学停止。}写出能直接核验的停止证书；预算耗尽不等同于收敛。\par
\textbf{输出。}列出返回对象、实际计数、中心状态和用于复核的诊断。
\end{keypointbox}

\begin{AlgorithmContract}{
  input={训练语料及活动词表、主题数$K$、严格正初始$\alpha,\varphi$、局部/外层双预算和容差、可复现启动表及Newton回溯预算},
  output={各启动最后合法状态与轨迹、同口径最佳$\widehat\alpha,\widehat\varphi,\gamma,\eta$或未定义标志、实际启动/外层/文档/主题/线搜索计数、规范状态与最终诊断},
  preconditions={语料词索引和维数合法，活动词表每项至少出现一次，$K,R_{\rm start}\ge1$，初值在整个活动词表上严格正归一，全部预算为非负整数且容差/并列规则冻结},
  initialization={清空$\mathtt{record}[1{:}R_{\rm start}]$并置总计数为$0$；每个记录独立保存$\mathtt{feasible}$、$\mathtt{status}$、最后合法参数、最终ELBO和诊断，启动内临时标志不得在跨启动选择中使用},
  loop={按启动、外层轮、文档、主题和Newton试探的命名循环推进；启动内失败先设置局部标志并带标签退出，退出前必须原子写入该启动记录},
  update={完整局部推断、临时主题矩阵、临时$\alpha$和新ELBO全部通过归一化/单调门禁后，才原子提交一次外层轮},
  domain={$\alpha$严格正，$\varphi$在活动词表上逐项严格正且逐主题归一，局部责任度和$\gamma$有限严格正，已提交ELBO有限且不异常下降},
  failure={词表/支持契约失败返回$\mathtt{invalid\_input}(\mathtt{support\_failure})$；其他全局输入非法返回$\mathtt{invalid\_input}$；启动源失败记录$\mathtt{random\_source\_failure}$；局部预算耗尽传播$\mathtt{budget\_stop}$；空主题、支持塌缩或非有限更新记录$\mathtt{numerical\_failure}$；Newton预算内无正参数上升步记录$\mathtt{line\_search\_failed}$，均保留上一合法启动状态},
  stop={每启动满足外层ELBO/参数双证书、外层预算耗尽或首次启动内失败；全部启动记录完成后按冻结口径选择},
  budget={至多$R_{\rm start}T_{\rm out}$个原子外层轮；每轮至多$M T_{\rm loc}$个局部轮、$K$个主题更新与$J_{\rm ls}$个Newton候选},
  status={$\mathtt{converged}$、$\mathtt{budget\_stop}$、$\mathtt{invalid\_input}$、$\mathtt{numerical\_failure}$、$\mathtt{line\_search\_failed}$、$\mathtt{random\_source\_failure}$},
  iterations={$s_{\rm done}$计实际记录启动，$r_{\rm done}$计原子外层轮，$m_{\rm done},k_{\rm done},j_{\rm done}$分别计完成文档、主题与Newton试探},
  diagnostics={$\mathtt{record}[s]$逐启动保存ELBO/双停止量、可行性、局部状态、空主题、Newton余量、失败$(s,r,m,k,j)$与命名循环阶段；跨启动选择只读记录数组并报告候选集合},
  complexity={最坏外层代价按$R_{\rm start}T_{\rm out}$乘以全部局部推断、稀疏词元责任度、主题归一和实际Newton试探成本累计}
}
% v2.6 layout: former forced clear removed; natural pagination.
\paragraph{稳健实现说明：多启动与受保护Newton。}
\begin{algorithm}[tbp]
\caption{点参数LDA变分EM（多启动）}\label{alg:V5-C06-vem}
\KwIn{语料$\mathcal D$；$K$；严格正初始$\boldsymbol\alpha,\boldsymbol\varphi$；
$\varepsilon_{\mathrm{loc}},\varepsilon_{\mathrm{out}},T_{\mathrm{loc}},T_{\mathrm{out}}$；
$R_{\mathrm{start}}$个可复现启动；Newton回溯上限$J_{\mathrm{ls}}$}
\KwOut{各启动最后合法状态与最佳参数；$s_{\rm done},r_{\rm done},m_{\rm done},k_{\rm done},j_{\rm done}$；$\mathtt{status}$与诊断}
冻结训练活动词表并核验其中每个词至少出现一次；再核验词索引、维数、$\alpha$正性、$\varphi$在整个活动词表上的严格正性、容差、预算与启动表。支持失败返回空结果、计数$0$、$\mathtt{invalid\_input}$及$\mathtt{support\_failure}$位置诊断，其余输入失败返回字段诊断\;
清空$\mathtt{record}[1{:}R_{\mathrm{start}}]$并置$s_{\rm done}=r_{\rm done}=m_{\rm done}=k_{\rm done}=j_{\rm done}=0$\;
\For{标记启动循环$\mathtt{start\_loop}$：$s\leftarrow1$ \KwTo $R_{\mathrm{start}}$}{
  置局部$\mathtt{startup\_ok}\leftarrow\mathtt{true}$并初始化该启动临时状态；置$\mathtt{record}[s].\mathtt{feasible}\leftarrow\mathtt{false}$、$\mathtt{record}[s].\mathtt{status}\leftarrow\mathtt{invalid\_input}$\;
  若启动随机源失败，置失败状态$\mathtt{random\_source\_failure}$和$\mathtt{startup\_ok}\leftarrow\mathtt{false}$\;
  若启动状态合法则计算基线ELBO；非有限时置失败状态$\mathtt{numerical\_failure}$和$\mathtt{startup\_ok}\leftarrow\mathtt{false}$\;
  \If{$\mathtt{startup\_ok}=\mathtt{true}$}{
    \For{标记外层循环$\mathtt{outer\_loop}$：$r\leftarrow1$ \KwTo $T_{\mathrm{out}}$}{
      \For{标记文档循环$\mathtt{document\_loop}$：$m\leftarrow1$ \KwTo $M$}{
        按式\eqref{eq:V5-C06-eta-update}--\eqref{eq:V5-C06-gamma-update}及局部实现条带执行单文档更新并增加$m_{\rm done}$；保存局部状态与诊断\;
        若局部未收敛，设置相应$\mathtt{budget\_stop}$或$\mathtt{numerical\_failure}$，令$\mathtt{startup\_ok}\leftarrow\mathtt{false}$并\textbf{break }$\mathtt{document\_loop}$\;
      }
      若$\mathtt{startup\_ok}=\mathtt{false}$，则\textbf{break }$\mathtt{outer\_loop}$；不得执行M步\;
      汇总临时$C_{kv}^+$\;
      \For{标记主题循环$\mathtt{topic\_loop}$：$k\leftarrow1$ \KwTo $K$}{
        若该行任一$C_{kv}^+$非有限，记录$(s,r,k,v)$与$\mathtt{numerical\_failure}(\mathtt{nonfinite\_count})$，令$\mathtt{startup\_ok}\leftarrow\mathtt{false}$并\textbf{break }$\mathtt{topic\_loop}$\;
        计算$C_{k\cdot}^+$；若$C_{k\cdot}^+\le0$，记录$(s,r,k)$与$\mathtt{numerical\_failure}(\mathtt{empty\_topic})$，令$\mathtt{startup\_ok}\leftarrow\mathtt{false}$并\textbf{break }$\mathtt{topic\_loop}$\;
        若该行任一$C_{kv}^+\le0$，记录$(s,r,k,v)$与$\mathtt{numerical\_failure}(\mathtt{support\_collapse})$，令$\mathtt{startup\_ok}\leftarrow\mathtt{false}$并\textbf{break }$\mathtt{topic\_loop}$\;
        对每个$v$在临时行形成严格正$\varphi_{kv}^+=C_{kv}^+/C_{k\cdot}^+$，核验整行正性与归一化后增加$k_{\rm done}$\;
      }
      若$\mathtt{startup\_ok}=\mathtt{false}$，则\textbf{break }$\mathtt{outer\_loop}$\;
      在至多$J_{\rm ls}$个候选内形成严格正$\alpha^+$并逐次增加$j_{\rm done}$；若无合格步，置$\mathtt{line\_search\_failed}$与$\mathtt{startup\_ok}=\mathtt{false}$并\textbf{break }$\mathtt{outer\_loop}$\;
      用完整临时候选计算ELBO；若非有限或异常下降，置$\mathtt{numerical\_failure}$与$\mathtt{startup\_ok}=\mathtt{false}$并\textbf{break }$\mathtt{outer\_loop}$\;
      原子提交$(\alpha,\varphi,\gamma,\eta,\mathrm{ELBO})$整轮候选并增加$r_{\rm done}$\;
      若外层双证书满足，置启动状态$\mathtt{converged}$并\textbf{break }$\mathtt{outer\_loop}$\;
    }
    若$\mathtt{startup\_ok}=\mathtt{true}$且未收敛，置启动状态$\mathtt{budget\_stop}$\;
  }
  原子写入$\mathtt{record}[s].\mathtt{feasible}\leftarrow\mathtt{startup\_ok}$、$\mathtt{record}[s].\mathtt{status}$、最后合法参数、最终ELBO、计数和诊断，并增加$s_{\rm done}$\;
  \If{$\mathtt{startup\_ok}=\mathtt{false}$}{在此唯一位置\textbf{continue }$\mathtt{start\_loop}$\;}
}
令$\mathcal S_{\rm acc}=\{s:\mathtt{record}[s].\mathtt{feasible}=\mathtt{true},\ \mathtt{record}[s].\mathtt{status}\in\{\mathtt{converged},\mathtt{budget\_stop}\}\}$；跨启动阶段只读$\mathtt{record}$，在评价口径相同的$\mathcal S_{\rm acc}$中按记录的最终ELBO和冻结并列规则选择\;
若$\mathcal S_{\rm acc}=\varnothing$，按全部$\mathtt{record}[s].\mathtt{status}$的确定性优先级返回$\mathtt{random\_source\_failure}$、$\mathtt{line\_search\_failed}$、$\mathtt{budget\_stop}$或$\mathtt{numerical\_failure}$及全部失败位置\;
返回最佳最后合法参数、所选启动规范状态、实际计数和跨启动最终诊断\;
\end{algorithm}
\end{AlgorithmContract}
\FloatBarrier
\SLRobustImplementationStrip{核心五步是逐文档局部E步、汇总期望计数、归一化M步、ELBO验收、在同一评价口径内选择多启动；36行故障树属于稳健实现细节}
\endgroup

\phantomsection\label{struct:V5-C06-CH28}
\textbf{适用条件。}LDA适合无监督探索中长文本集合、词袋近似尚可、主题可由共现词解释且主题数可通过验证选择的场景。
短文本、强顺序语义、动态主题、监督预测或词义高度上下文化时，需要聚合文本、序列主题模型、动态/监督扩展或上下文表示；不能只调$K$掩盖模型错配。

\phantomsection\label{struct:V5-C06-CH29}
\textbf{局限性。}主题标签置换不可识别；似然非凸导致局部最优；Gibbs可能混合缓慢，平均场常低估后验相关与方差；固定$K$、词袋假设和单一词义限制表达；高频停用词、极稀有词和语料漂移会扭曲主题。主题“可解释”仍需人工和稳定性检验，不能由困惑度单独保证。

\phantomsection\label{struct:V5-C06-CH30}
\begin{pitfallbox}
Gibbs计算前未删当前词元；把$n_{mk}$和$n_{kv}$方向写反；遗漏$\boldsymbol\beta$；
把完整Bayes Gibbs与点参数变分EM写成同一个后验；把未归一化期望计数当$\boldsymbol\varphi$；
在$\eta$更新中漏掉digamma或归一化；把ELBO当真实对数证据；
Newton步产生非正$\alpha$仍继续；用训练词同时推断文档主题并评价同一词的预测概率；只跑一条链或一次初始化便宣布稳定。
\end{pitfallbox}

\phantomsection\label{struct:V5-C06-CH31}
% KNOWLEDGE-PLAN: KN-V5-C35-COMPARISON-001 L1
\paragraph{读前自检：易混淆概念对比：稳健实现说明：多启动与受保护Newton。}围绕稳健实现说明：多启动与受保护Newton，先采用条件“对比表首个数据行是“$\boldsymbol\theta_m$与$\boldsymbol\varphi_k$｜前者是文档在主题上的分布，维数$K$｜后者是主题在词表上的分布，维数$V$””，再按列复述“$\boldsymbol\theta_m$与$\boldsymbol\varphi_k$｜前者是文档在主题上的分布，维数$K$｜后者是主题在词表上的分布，维数$V$”中的对象、假设与结论；最后验证各列均对应首行对象“$\boldsymbol\theta_m$与$\boldsymbol\varphi_k$”，且未与表中下一行互换。\par

\begin{comparisonbox}
\small
\begin{tabularx}{\textwidth}{@{}p{0.18\textwidth}XX@{}}
\toprule
对象 & 第一项 & 第二项 \\
\midrule
$\boldsymbol\theta_m$与$\boldsymbol\varphi_k$ & 前者是文档在主题上的分布，维数$K$ & 后者是主题在词表上的分布，维数$V$ \\
$\boldsymbol\alpha$与$\boldsymbol\beta$ & 前者平滑文档--主题计数 & 后者平滑主题--词计数 \\
$\boldsymbol\eta_{mn}$与$\boldsymbol\varphi_k$ & 前者是某词元属于主题的变分责任度 & 后者是给定主题生成词的模型概率 \\
折叠与非折叠 & 折叠法积分消去共轭概率向量 & 非折叠法显式抽样或近似它们 \\
Gibbs与变分 & 前者输出相关后验样本并需链诊断 & 后者输出确定性近似并需ELBO与多启动诊断 \\
困惑度与解释性 & 困惑度衡量留出预测 & 解释性关注语义连贯、稳定与人工可读性 \\
\bottomrule
\end{tabularx}
\end{comparisonbox}

评估时，把训练集之外的词元记为$\mathcal H$，用训练部分或文档已观察部分推断
$\widehat{\boldsymbol\theta}_m$，再计算
\begin{equation}
\operatorname{perplexity}(\mathcal H)
=\exp\!\left\{-\frac{1}{|\mathcal H|}
\sum_{(m,n)\in\mathcal H}
\log\sum_{k=1}^{K}\widehat\theta_{mk}\widehat\varphi_{k,w_{mn}}
\right\}.
\label{eq:V5-C06-perplexity}
\end{equation}
必须固定切分、词表和未知词处理；文档完成式评价不能让留出词参与
$\widehat{\boldsymbol\theta}_m$的推断。对每个留出词还须在取对数前核验
$w_{mn}\in\mathcal V_{\mathrm{tr}}$（或已经映射到训练中出现过的$\mathtt{UNK}$）以及
$\sum_k\widehat\theta_{mk}\widehat\varphi_{k,w_{mn}}>0$；任一条件失败就返回
$\mathtt{invalid\_input}(\mathtt{support\_failure},m,n,w_{mn})$，困惑度记为未定义，不能把$\log0$送入浮点计算。

\paragraph{公平比较原则。}
比较两条路线时，应使用同一训练/验证/测试切分，并只在训练数据上确定词表与预处理规则；固定$K$、初始化原则和计算预算。Gibbs法报告链间一致性、自相关或有效样本量，变分法报告ELBO轨迹与不同初值的离散程度；两者使用同一困惑度、主题稳定性和语义连贯度口径。
还必须把“完整Bayes LDA”和“点参数LDA变体”登记为不同模型配置；否则二者结果差异同时包含模型目标与推断方法变化，不能归因于推断算法。若词表、遮罩、预算或主题匹配规则不同，同样不能作单因素归因。


\Needspace{4\baselineskip}
\SLDirectSection{例题与练习}{sec:V5-C06-S06}
\begin{SLRunningExample}{RUN-DOC-TERM-01}{V5-C06-lda}{贯穿例：文档—词项矩阵小例}{把同一非负整数词项计数矩阵展开为有序词元，并为主题比例与主题词分布加入Dirichlet随机层。}
沿用词为行、文档为列的计数矩阵
\[
\boldsymbol X=\begin{bmatrix}2&0&1\\1&2&0\\0&1&3\end{bmatrix}.
\]
矩阵元素均为非负整数，三列总数分别为$3,3,4$且没有空文档，因而既符合上一站的词项计数约定，也能合法展开为LDA的有序观测词元。本章不把它只写成$WH$，而是分别用$\boldsymbol\theta_m$与$\boldsymbol\varphi_k$生成主题指派和单词；矩阵总计数等于词元数$I=10$。上一站见\SLRunningExampleRef{RUN-DOC-TERM-01}{V4-C06-plsa}{第4册第6章的PLSA概率分解}。
\end{SLRunningExample}

\phantomsection\label{struct:V5-C06-CH32}
\Needspace{6\baselineskip}
\begin{example}[一次折叠Gibbs更新]\label{exm:V5-C06-gibbs-step}
取本章小型例子的计数与超参数。求两个候选主题的满条件概率，并说明旧主题计数何时删除、何时恢复。
\end{example}
\SLExampleSolutionHeading{exm:V5-C06-gibbs-step}
\begin{solution}
\SLReadTranslation{题目要求完成“一次折叠Gibbs更新”：先写清目标分布、提议或满条件与支持，再构造转移/估计量并核验不变性或归一化。}
\SolGiven 目标分布、提议核或满条件在题面支持上有定义；所有条件分母和归一化常数为正，状态空间与随机变量取值域保持一致。
\SLMethodTrigger{先声明目标分布、提议/满条件与支持，再逐步构造转移或估计量并核验归一化。}
\SolDerive
\textbf{计算。}
删去当前词元后按式\eqref{eq:V5-C06-gibbs-conditional}得
\[
r_1=\frac3{26},\qquad r_2=\frac5{44},\qquad
r_1+r_2=\frac{131}{572}.
\]
归一化后
\[
\Pr(z_i=1\mid\boldsymbol z_{-i},\boldsymbol w)
=\frac{r_1}{r_1+r_2}=\frac{66}{131},\qquad
\Pr(z_i=2\mid\boldsymbol z_{-i},\boldsymbol w)=\frac{65}{131}.
\]
必须在计算前从旧主题对应的两张表各减一；抽到新主题后再加一。
\textbf{独立核验。}$66/131+65/131=1$；若把抽得的新主题再删去并把旧主题加回，两张计数表都精确恢复到更新前状态，说明“减--算--抽--加”保持计数守恒。

\textbf{结论。}本次候选概率为$(66/131,65/131)$，且删除与恢复必须成对执行；若未先减一或减一后不恢复，就不再是同一个Gibbs核。
\end{solution}
\Needspace{6\baselineskip}
\begin{example}[一次平均场责任度更新]\label{exm:V5-C06-vi-step}
给定$K=2$、$\boldsymbol\gamma=(3,2)$以及
$(\varphi_{1v},\varphi_{2v})=(0.6,0.2)$，计算$\boldsymbol\eta_n$。
\end{example}
\SLExampleSolutionHeading{exm:V5-C06-vi-step}
\begin{solution}
\SLReadTranslation{题目要求完成“一次平均场责任度更新”：依据题面概率模型求出目标量，并解释条件化、归一化或边缘化。}
\SolGiven 题面概率、权重或密度均处于合法范围；条件分母/归一化常数应为正，支持集与随机变量取值域保持一致。
\SLMethodTrigger{先写支持集、条件分母或归一化常数，再代入概率公式并做概率和/边缘核验。}
\SolDerive
\textbf{计算。}
整数digamma差给出
\[
\psi(3)-\psi(5)=-\frac7{12},\qquad
\psi(2)-\psi(5)=-\frac{13}{12},
\]
故未归一化权重和责任度依次为
\[
\boldsymbol u
=\bigl(0.6e^{-7/12},\,0.2e^{-13/12}\bigr)
\approx(0.3348,0.0677),\qquad
\boldsymbol\eta_n
=\frac{\boldsymbol u}{u_1+u_2}
\approx(0.832,0.168).
\]
\textbf{独立核验。}两个主题词概率都严格为正，且
$u_1/u_2=3e^{1/2}\approx4.946$与$0.832/0.168$在舍入误差内一致；同时$0.832+0.168=1$。

\textbf{结论。}责任度约为$(0.832,0.168)$；归一化和严格正支持都是坐标更新的定义条件，不是事后的显示修正。
\end{solution}
\Needspace{6\baselineskip}
\begin{example}[留出困惑度]\label{exm:V5-C06-perplexity}
两个留出词的预测概率分别为$0.25$和$0.5$。求每词困惑度；若错误地把概率改成训练后重估得到的$0.5$和$0.8$，说明为什么数值虽更小却不可比较。
\end{example}
\SLExampleSolutionHeading{exm:V5-C06-perplexity}
\begin{solution}
\SLReadTranslation{题目要求完成“留出困惑度”：依据题面概率模型求出目标量，并解释条件化、归一化或边缘化。}
\SolGiven 题面概率、权重或密度均处于合法范围；条件分母/归一化常数应为正，支持集与随机变量取值域保持一致。
\SLMethodTrigger{先写支持集、条件分母或归一化常数，再代入概率公式并做概率和/边缘核验。}
\SolDerive
\textbf{计算。}
\[
\operatorname{perplexity}
=\exp\{-\tfrac12(\log0.25+\log0.5)\}
=\frac1{\sqrt{0.125}}=2\sqrt2\approx2.828.
\]
若让留出词参与重估，则表面数值变成
\[
\operatorname{perplexity}_{\rm leak}
=\exp\{-\tfrac12(\log0.5+\log0.8)\}
=\frac1{\sqrt{0.4}}\approx1.581.
\]
\textbf{独立核验。}合法值也可由预测概率几何平均复算：
$\operatorname{perplexity}=1/\sqrt{0.25\cdot0.5}=2\sqrt2$；两个概率均严格为正，因此没有支持失败。

\textbf{结论。}合法留出困惑度约为$2.828$。后一个$1.581$使用留出词重估参数，信息泄漏使概率被人为抬高，所得数值不再是同一预测任务的留出指标。
\end{solution}

\phantomsection\label{struct:V5-C06-CH33}
\Needspace{4\baselineskip}
\phantomsection\label{struct:V5-C06-CH34}
\begin{chapterexercisebox}
\phantomsection\label{struct:V5-C06-CH35}
本章共16道练习。每道题干后立即给出同号解析；长解析可自然跨页，续页标题保留题号。
\end{chapterexercisebox}

\begin{SLChapterExercisePairSection}

\SLSourceBookExercises

\Needspace{6\baselineskip}
% EXERCISE-SOURCE: original_book; source_id=EXR-20-0002; source_number=20.2
% SOLUTION-SOURCE: original_book; source_id=EXR-20-0002; source_number=20.2
\begin{exercise}\label{ex:V5-C06-S01-01}
对第四册潜在语义分析章使用的固定小语料建立$K=2$的LDA分析方案：给出词表、生成变量、一次可复现初始化、至少一次推断更新以及文档主题和主题关键词的输出格式。若不实际运行随机程序，哪些结论只能写成“计算方案”而不能虚构成数值结果？
\end{exercise}
\ifSLFullEdition
\phantomsection\label{sol:V5-C06-S01-01}
\begin{chapterexercisesolution}{ex:V5-C06-S01-01}
\nopagebreak[4]
沿用第四册固定的词袋矩阵，按下列顺序编码词表：
（\texttt{book}, \texttt{dads}, \texttt{dummies}, \texttt{estate},
\texttt{guide}, \texttt{investing}, \texttt{market}, \texttt{real},
\texttt{rich}, \texttt{stock}, \texttt{value}）。九篇文档依次为
$T1=$（\texttt{guide}, \texttt{investing}, \texttt{market}, \texttt{stock}），
$T2=$（\texttt{dummies}, \texttt{investing}），
$T3=$（\texttt{book}, \texttt{investing}, \texttt{market}, \texttt{stock}），
$T4=$（\texttt{book}, \texttt{investing}, \texttt{value}），
$T5=$（\texttt{investing}, \texttt{value}），
$T6=$（\texttt{dads}, \texttt{guide}, \texttt{investing}, \texttt{rich}, \texttt{rich}），
$T7=$（\texttt{estate}, \texttt{investing}, \texttt{real}），
$T8=$（\texttt{dummies}, \texttt{investing}, \texttt{stock}），
$T9=$（\texttt{dads}, \texttt{estate}, \texttt{investing}, \texttt{real}, \texttt{rich}），
共31个词元。

取$K=2$、$\boldsymbol\alpha=(1,1)$、$\boldsymbol\beta=\boldsymbol1_{11}$；
每篇文档内部按上述词表顺序展开词元，并令
$z_{mn}^{(0)}=1$当且仅当$m+n$为偶数，否则令其为2。这给出不依赖随机数的可复现初始化。
以$T1$首词\texttt{guide}为一次手算更新：删去它后，
$(n_{11}^{-i},n_{12}^{-i})=(1,2)$、
$(n_{1,\mathrm{guide}}^{-i},n_{2,\mathrm{guide}}^{-i})=(1,0)$、
$(n_{1\cdot}^{-i},n_{2\cdot}^{-i})=(14,16)$，故两个未归一化权重为
\[
r_1=\frac{1+1}{3+2}\frac{1+1}{14+11}=\frac4{125},\qquad
r_2=\frac{2+1}{3+2}\frac{0+1}{16+11}=\frac1{45}.
\]
归一化链为
\[
r_1+r_2=\frac{61}{1125},\qquad
\Pr(z_i=1\mid-)=\frac{36}{61},\qquad
\Pr(z_i=2\mid-)=\frac{25}{61}.
\]
这一步核验的是“减--算”；实际“抽--加”
还须记录伪随机生成器、实现版本与种子，不能把未执行的随机抽样写成既成结果。

最终输出应固定为两张表：$T_m\mid\widehat\theta_{m1},\widehat\theta_{m2}$，以及
$k\mid$按$\widehat\varphi_{kv}$降序排列的词与概率；另附迭代数、留样规则、种子和诊断量。
其中每个留样状态的归一化估计量是
\[
\widehat\theta_{mk}
=\frac{n_{mk}+\alpha_k}{N_m+\alpha_0},\qquad
\widehat\varphi_{kv}
=\frac{n_{kv}+\beta_v}{n_{k\cdot}+\beta_0},
\]
多次留样时再对这些估计量作样本平均。
变分方案也须记录初始$\boldsymbol\varphi$，按算法\ref{alg:V5-C06-vem}完成局部更新和一次归一化期望计数M步。
若没有实际运行程序，只能报告上述数据编码、公式、手算条件更新、输出表结构和验收条件；
不能声称某词最终属于某主题、某文档比例为何值，也不能虚构困惑度或收敛轮数。
\end{chapterexercisesolution}
\fi

\Needspace{6\baselineskip}
% EXERCISE-SOURCE: original_book; source_id=EXR-20-0003; source_number=20.3
% SOLUTION-SOURCE: original_book; source_id=EXR-20-0003; source_number=20.3
\begin{exercise}\label{ex:V5-C06-S03-01}
分别指出折叠Gibbs与变分EM中Dirichlet分布出现的位置，并解释共轭性、平滑和期望对数各自为何重要。
\end{exercise}
\ifSLFullEdition
\phantomsection\label{sol:V5-C06-S03-01}
\begin{chapterexercisesolution}{ex:V5-C06-S03-01}
\nopagebreak[4]
Gibbs中$\boldsymbol\theta_m\sim\operatorname{Dir}(\boldsymbol\alpha)$、
$\boldsymbol\varphi_k\sim\operatorname{Dir}(\boldsymbol\beta)$与类别似然共轭，使二者可积分消去，留下Beta函数比值和式\eqref{eq:V5-C06-gibbs-conditional}；
$\boldsymbol\alpha,\boldsymbol\beta$又以伪计数平滑零计数。
\[
\int p(\boldsymbol\theta_m\mid\boldsymbol\alpha)
\prod_k\theta_{mk}^{n_{mk}}\,\mathrm d\boldsymbol\theta_m
=\frac{B(\boldsymbol n_m+\boldsymbol\alpha)}{B(\boldsymbol\alpha)},\qquad
\int p(\boldsymbol\varphi_k\mid\boldsymbol\beta)
\prod_v\varphi_{kv}^{n_{kv}}\,\mathrm d\boldsymbol\varphi_k
=\frac{B(\boldsymbol n_k+\boldsymbol\beta)}{B(\boldsymbol\beta)}.
\]

变分法中$q(\boldsymbol\theta_m)=\operatorname{Dir}(\boldsymbol\gamma_m)$，其
\[
\mathbb E_q\log\theta_{mk}
=\psi(\gamma_{mk})-\psi(\gamma_{m0}),\qquad
\eta_{mnk}\propto
\varphi_{k,w_{mn}}\exp\{\mathbb E_q\log\theta_{mk}\},
\]
进入责任度更新；模型先验仍通过
\[
\gamma_{mk}=\alpha_k+\sum_n\eta_{mnk},\qquad
\sum_k\eta_{mnk}=1
\]
提供平滑。
因此共轭性保证解析族闭合，平滑避免零概率，期望对数则是平均场坐标更新的自然参数，三者作用不同。
\end{chapterexercisesolution}
\fi

\Needspace{6\baselineskip}
% EXERCISE-SOURCE: original_book; source_id=EXR-20-0004; source_number=20.4
% SOLUTION-SOURCE: original_book; source_id=EXR-20-0004; source_number=20.4
\begin{exercise}\label{ex:V5-C06-S05-01}
以$I,K,V,M,T_{\mathrm G},T_{\mathrm E},T_{\mathrm M}$为规模参数，给出稠密折叠Gibbs和批量变分EM的训练时间、预测时间与空间复杂度，并说明每个界的实现假设。
\end{exercise}
\ifSLFullEdition
\phantomsection\label{sol:V5-C06-S05-01}
\begin{chapterexercisesolution}{ex:V5-C06-S05-01}
\nopagebreak[4]
令$T_{\mathrm G}$为Gibbs总扫描数，则稠密候选计算是
$O(T_{\mathrm G}IK)$。更完整地，
\[
\text{Gibbs训练}=O(T_{\mathrm G}IK),\qquad
\text{新文档折入}=O(T_{\mathrm E}N_mK),\qquad
\text{工作空间}=O(KV+MK+I).
\]
若$T_{\mathrm S}$个完整指派样本全部驻留内存，还须加$O(T_{\mathrm S}I)$，流式落盘或在线汇总则不加。
令变分EM外层为$T_{\mathrm M}$、每个E步平均$T_{\mathrm E}$轮，则训练为
\[
\text{变分训练}
=O(T_{\mathrm M}T_{\mathrm E}IK+T_{\mathrm M}KV),\qquad
\text{新文档推断}=O(T_{\mathrm E}N_mK).
\]
批量保存责任度与流式处理的空间界分别为
\[
O(IK+KV+MK),\qquad
O(K\max_mN_m+KV+MK).
\]
这里还假设每个外层轮次只做固定上限个Newton候选，并利用Hessian的对角加秩一结构以$O(K)$求解；若用普通稠密求解，训练界另加$O(T_{\mathrm M}K^3)$。
这些界假设稠密遍历$K$个主题、固定精度特殊函数和$O(1)$数组访问；稀疏/并行实现需另列同步与索引成本。
\end{chapterexercisesolution}
\fi

\Needspace{6\baselineskip}
% EXERCISE-SOURCE: original_book; source_id=EXR-20-0005; source_number=20.5
% SOLUTION-SOURCE: original_book; source_id=EXR-20-0005; source_number=20.5
\begin{exercise}\label{ex:V5-C06-S04-01}
证明精确块更新下变分EM的ELBO单调不减；说明为什么这不自动证明真实证据在每个受限E步后等于ELBO，也不保证参数收敛到全局最优。
\end{exercise}
\ifSLFullEdition
\phantomsection\label{sol:V5-C06-S04-01}
\begin{chapterexercisesolution}{ex:V5-C06-S04-01}
\nopagebreak[4]
固定模型参数的精确块E更新不会降低ELBO；固定新$q$的精确块M更新也不会降低ELBO，故
\[
\mathcal L(q^{t},\Theta^{t})
\le\mathcal L(q^{t+1},\Theta^{t})
\le\mathcal L(q^{t+1},\Theta^{t+1}).
\]
ELBO有上界时其数值收敛。由式\eqref{eq:V5-C06-elbo-identity}，
\[
\log p(\boldsymbol w\mid\Theta)
=\mathcal L(q,\Theta)
+\KL\!\left(q(\boldsymbol h)\,\middle\|\,p(\boldsymbol h\mid\boldsymbol w,\Theta)\right),
\]
ELBO等于真实证据还要求KL间隙为零；受限平均场通常不满足。
非凸目标、标签置换和边界驻点又说明ELBO值收敛不等于参数唯一，更不等于全局最优。
\end{chapterexercisesolution}
\fi

\SLLectureSupplementExercises

\Needspace{6\baselineskip}
% EXERCISE-SOURCE: lecture_supplement
% SOLUTION-SOURCE: lecture_supplement
\begin{exercise}\label{ex:V5-C06-S02-01}
写出$M=3,K=4,V=10$时$\boldsymbol\theta,\boldsymbol\varphi,\boldsymbol z,\boldsymbol w$的维数与取值域，并把式\eqref{eq:V5-C06-joint}按板块逐项核对。
\end{exercise}
\ifSLFullEdition
\phantomsection\label{sol:V5-C06-S02-01}
\begin{chapterexercisesolution}{ex:V5-C06-S02-01}
\nopagebreak[4]
\[
\boldsymbol\theta=(\theta_{mk})\in(\Delta^3)^3
\quad(3\times4),\qquad
\boldsymbol\varphi=(\varphi_{kv})\in(\Delta^9)^4
\quad(4\times10).
\]
两组离散变量按文档长度形成锯齿形序列：
\[
z_{mn}\in\{1,2,3,4\},\qquad
w_{mn}\in\{1,\ldots,10\},\qquad
n=1,\ldots,N_m.
\]
式\eqref{eq:V5-C06-joint}逐板块展开为
\[
\prod_{k=1}^{4}p(\boldsymbol\varphi_k\mid\boldsymbol\beta)
\prod_{m=1}^{3}p(\boldsymbol\theta_m\mid\boldsymbol\alpha)
\prod_{m=1}^{3}\prod_{n=1}^{N_m}
p(z_{mn}\mid\boldsymbol\theta_m)
p(w_{mn}\mid z_{mn},\boldsymbol\varphi).
\]
因此有4个主题先验因子、3个文档先验因子和$\sum_mN_m$对“主题指派×单词发射”因子，恰与三个板块重复次数一致。
\end{chapterexercisesolution}
\fi

\Needspace{6\baselineskip}
% EXERCISE-SOURCE: lecture_supplement
% SOLUTION-SOURCE: lecture_supplement
\begin{exercise}\label{ex:V5-C06-S02-02}
从概率图判断：$w_{mn}$在给定哪些变量时与其他词位置条件独立？边缘化$\boldsymbol\theta_m$后，同文档两个词是否仍独立？给出理由。
\end{exercise}
\ifSLFullEdition
\phantomsection\label{sol:V5-C06-S02-02}
\begin{chapterexercisesolution}{ex:V5-C06-S02-02}
\nopagebreak[4]
给定$z_{mn}$和全部$\varphi$时，$w_{mn}$与其他位置的词和主题条件独立；若不显式给$z_{mn}$，给定$\theta_m,\varphi$后积分该位置主题，词位置仍条件独立。
\[
p(\boldsymbol w_m\mid\boldsymbol z_m,\boldsymbol\varphi)
=\prod_{n=1}^{N_m}\varphi_{z_{mn},w_{mn}},\qquad
p(\boldsymbol w_m\mid\boldsymbol\theta_m,\boldsymbol\varphi)
=\prod_{n=1}^{N_m}\sum_k\theta_{mk}\varphi_{k,w_{mn}}.
\]
边缘化共享的$\theta_m$后，同文档位置通常相关，因为观察一个词会改变对$\theta_m$的后验，从而改变另一词的预测分布。
\[
p(w_{m1},w_{m2}\mid\boldsymbol\varphi,\boldsymbol\alpha)
=\int
\prod_{j=1}^{2}\left(\sum_k\theta_{mk}\varphi_{k,w_{mj}}\right)
p(\boldsymbol\theta_m\mid\boldsymbol\alpha)\,
\mathrm d\boldsymbol\theta_m
\ne
\prod_{j=1}^{2}p(w_{mj}\mid\boldsymbol\varphi,\boldsymbol\alpha)
\]
一般成立，右侧不等是共享随机混合比例造成的。
\end{chapterexercisesolution}
\fi

\Needspace{6\baselineskip}
% EXERCISE-SOURCE: lecture_supplement
% SOLUTION-SOURCE: lecture_supplement
\begin{exercise}\label{ex:V5-C06-S02-03}
证明给定$\boldsymbol\theta_m,\boldsymbol\varphi$时词序列交换；再构造一个共享随机主题比例的二词例子，说明交换不等于边缘独立。
\end{exercise}
\ifSLFullEdition
\phantomsection\label{sol:V5-C06-S02-03}
\begin{chapterexercisesolution}{ex:V5-C06-S02-03}
\nopagebreak[4]
给定参数时联合概率是同一组位置因子的乘积，置换不改变乘积。
\[
p(w_{m1:N_m}\mid\boldsymbol\theta_m,\boldsymbol\varphi)
=\prod_{n=1}^{N_m}\sum_k\theta_{mk}\varphi_{k,w_{mn}}
=p(w_{m,\pi(1):\pi(N_m)}\mid\boldsymbol\theta_m,\boldsymbol\varphi)
\]
对任意置换$\pi$都成立。
反例取随机$\Theta\sim\operatorname{Beta}(1,1)$，给定$\Theta$后两个二元词独立且成功概率为$\Theta$。
\[
\Pr(W_1=1)=\Pr(W_2=1)=\mathbb E\Theta=\frac12,\qquad
\Pr(W_1=W_2=1)=\mathbb E\Theta^2=\frac13\ne\frac14.
\]
所以两词边缘不独立；但联合式对两位置对称，故仍交换。
\end{chapterexercisesolution}
\fi

\Needspace{6\baselineskip}
% EXERCISE-SOURCE: lecture_supplement
% SOLUTION-SOURCE: lecture_supplement
\begin{exercise}\label{ex:V5-C06-S03-02}
从两个Dirichlet积分完整推出式\eqref{eq:V5-C06-collapsed-joint}，写清每个Beta函数的向量维数和积分变量。
\end{exercise}
\ifSLFullEdition
\phantomsection\label{sol:V5-C06-S03-02}
\begin{chapterexercisesolution}{ex:V5-C06-S03-02}
\nopagebreak[4]
对每个主题$k$，在$V$分量概率向量上积分得
\[
\int_{\Delta^{V-1}}
\frac1{B(\boldsymbol\beta)}\prod_v\varphi_{kv}^{n_{kv}+\beta_v-1}
\,\mathrm d\boldsymbol\varphi_k
=\frac{B(\boldsymbol n_k+\boldsymbol\beta)}{B(\boldsymbol\beta)}.
\]
对每篇文档$m$，在$K$分量概率向量上积分得
\[
\int_{\Delta^{K-1}}
\frac1{B(\boldsymbol\alpha)}\prod_k\theta_{mk}^{n_{mk}+\alpha_k-1}
\,\mathrm d\boldsymbol\theta_m
=\frac{B(\boldsymbol n_m+\boldsymbol\alpha)}{B(\boldsymbol\alpha)}.
\]
分别对$K$个主题和$M$篇文档相乘即得式\eqref{eq:V5-C06-collapsed-joint}。
\end{chapterexercisesolution}
\fi

\Needspace{6\baselineskip}
% EXERCISE-SOURCE: lecture_supplement
% SOLUTION-SOURCE: lecture_supplement
\begin{exercise}\label{ex:V5-C06-S03-03}
对当前词元$i=(m,n)$，用Gamma函数比值从折叠联合分布推出式\eqref{eq:V5-C06-gibbs-conditional}；解释上标$^{-i}$为何必须同时作用于两类计数。
\end{exercise}
\ifSLFullEdition
\phantomsection\label{sol:V5-C06-S03-03}
\begin{chapterexercisesolution}{ex:V5-C06-S03-03}
\nopagebreak[4]
候选$k$相对留一状态只把$n_{kv}^{-i}$与$n_{mk}^{-i}$各加1。
Beta函数比值用$\Gamma(x+1)/\Gamma(x)=x$化为
\[
r_k=
\frac{n_{kv}^{-i}+\beta_v}{n_{k\cdot}^{-i}+\beta_0}
\frac{n_{mk}^{-i}+\alpha_k}{n_{m\cdot}^{-i}+\alpha_0}.
\]
满条件概率还须在主题上归一化：
\[
\Pr(z_i=k\mid\boldsymbol z_{-i},\boldsymbol w)
=\frac{r_k}{\sum_{\ell=1}^{K}r_\ell}.
\]
若只从一张表删旧词元，两张表不再描述同一个$\boldsymbol z_{-i}$，所得比值就不是任何一致联合分布的满条件概率。
\end{chapterexercisesolution}
\fi

\Needspace{6\baselineskip}
% EXERCISE-SOURCE: lecture_supplement
% SOLUTION-SOURCE: lecture_supplement
\begin{exercise}\label{ex:V5-C06-S03-04}
取$K=2,V=3$，$\alpha=(1,1)$、$\beta=(1,1,1)$，留一计数
$(n_{1v},n_{2v})=(2,1)$、$(n_{1\cdot},n_{2\cdot})=(5,4)$、
$(n_{m1},n_{m2})=(1,3)$。计算一次主题更新的归一化概率。
\end{exercise}
\ifSLFullEdition
\phantomsection\label{sol:V5-C06-S03-04}
\begin{chapterexercisesolution}{ex:V5-C06-S03-04}
\nopagebreak[4]
两个权重为
\[
r_1=\frac{2+1}{5+3}\frac{1+1}{1+3+2}
=\frac18,
\qquad
r_2=\frac{1+1}{4+3}\frac{3+1}{1+3+2}
=\frac4{21}.
\]
归一化链为
\[
r_1+r_2=\frac{53}{168},\qquad
p(z_i=1\mid-)=\frac{21}{53},\qquad
p(z_i=2\mid-)=\frac{32}{53}.
\]
\end{chapterexercisesolution}
\fi

\Needspace{6\baselineskip}
% EXERCISE-SOURCE: lecture_supplement
% SOLUTION-SOURCE: lecture_supplement
\begin{exercise}\label{ex:V5-C06-S03-05}
给定$\alpha=(2,1)$、$\beta=(1,1,1)$，某文档主题计数$(3,2)$，两个主题的词计数行分别为$(4,1,0)$和$(1,2,2)$。计算平滑后的$\widehat{\boldsymbol\theta}_m$与两行$\widehat{\boldsymbol\varphi}$，核验各行和为1。
\end{exercise}
\ifSLFullEdition
\phantomsection\label{sol:V5-C06-S03-05}
\begin{chapterexercisesolution}{ex:V5-C06-S03-05}
\nopagebreak[4]
文档后验伪计数为$(5,3)$，故
\[
\widehat{\boldsymbol\theta}_m
=\frac{(3,2)+(2,1)}{3+2+2+1}
=\left(\frac58,\frac38\right).
\]
第一主题词伪计数为$(5,2,1)$，第二主题为$(2,3,3)$，所以
\[
\widehat\varphi_1=(5/8,1/4,1/8),\qquad
\widehat\varphi_2=(1/4,3/8,3/8).
\]
归一化核验为
\[
\sum_k\widehat\theta_{mk}=1,\qquad
\sum_v\widehat\varphi_{1v}=1,\qquad
\sum_v\widehat\varphi_{2v}=1.
\]
\end{chapterexercisesolution}
\fi

\Needspace{6\baselineskip}
% EXERCISE-SOURCE: lecture_supplement
% SOLUTION-SOURCE: lecture_supplement
\begin{exercise}\label{ex:V5-C06-S04-02}
从KL定义证明式\eqref{eq:V5-C06-elbo-identity}，并分别说明$\mathcal Q$包含精确后验和不包含精确后验时最优KL间隙。
\end{exercise}
\ifSLFullEdition
\phantomsection\label{sol:V5-C06-S04-02}
\begin{chapterexercisesolution}{ex:V5-C06-S04-02}
\nopagebreak[4]
把$\log p(h\mid w)=\log p(w,h)-\log p(w)$代入
$\KL(q\|p)=\mathbb E_q\log q-\mathbb E_q\log p(h\mid w)$，得到
\[
\begin{aligned}
\KL\!\left(q(h)\,\middle\|\,p(h\mid w)\right)
&=\mathbb E_q\log q(h)-\mathbb E_q\log p(w,h)+\log p(w)\\
&=\log p(w)-\mathcal L(q),
\end{aligned}
\]
即
\[
\log p(w)=\mathcal L(q)
+\KL\!\left(q(h)\,\middle\|\,p(h\mid w)\right).
\]
若$\mathcal Q$含精确后验，最优$q$可使KL为0，ELBO达到证据；
对任一可行$q\ne p(h\mid w)$，KL严格为正。若$\mathcal Q$中的最小值实际取到且精确后验不属于$\mathcal Q$，
则最优间隙严格为正；若只有下确界，则该下确界也可能为0而不取到。
\end{chapterexercisesolution}
\fi

\Needspace{6\baselineskip}
% EXERCISE-SOURCE: lecture_supplement
% SOLUTION-SOURCE: lecture_supplement
\begin{exercise}\label{ex:V5-C06-S04-03}
利用拉格朗日乘子推出式\eqref{eq:V5-C06-eta-update}，必须写出$\sum_k\eta_{mnk}=1$约束和归一化常数。再按本章点参数模型的支持契约，分别判断“某主题对观察词概率为零”“留出词不在训练活动词表”“某主题期望总计数为零”三类输入应在何处停止。
\end{exercise}
\ifSLFullEdition
\phantomsection\label{sol:V5-C06-S04-03}
\begin{chapterexercisesolution}{ex:V5-C06-S04-03}
\nopagebreak[4]
与$\eta_{mn}$有关的目标为
\[
\sum_k\eta_{mnk}
\{\log\varphi_{k,w_{mn}}+\psi(\gamma_{mk})-\psi(\gamma_{m0})\}
-\sum_k\eta_{mnk}\log\eta_{mnk}
+\lambda(\sum_k\eta_{mnk}-1).
\]
记$a_k=\log\varphi_{k,w_{mn}}+\psi(\gamma_{mk})-\psi(\gamma_{m0})$，求偏导得
\[
\frac{\partial}{\partial\eta_{mnk}}
=a_k-\log\eta_{mnk}-1+\lambda=0,
\qquad
\eta_{mnk}=e^{\lambda-1}e^{a_k}.
\]
由$\sum_k\eta_{mnk}=1$确定归一化常数：
\[
u_k=e^{a_k},\qquad
Z_{mn}=\sum_{\ell=1}^{K}u_\ell,\qquad
\eta_{mnk}=\frac{u_k}{Z_{mn}},
\]
即式\eqref{eq:V5-C06-eta-update}。

\textbf{边界核验。}三类测试不能进入同一条浮点归一化链：
\textbf{（1）}若某个$\varphi_{k,w_{mn}}=0$，在局部更新取对数之前返回
\[
\mathtt{invalid\_input}(\mathtt{support\_failure}).
\]
\textbf{（2）}若留出词不属于$\mathcal V_{\rm tr}$且未映射到已训练
$\mathtt{UNK}$，在困惑度取对数之前返回同一
$\mathtt{support\_failure}$，困惑度未定义。
\textbf{（3）}若某主题$C_{k\cdot}=0$，在M步除法之前返回
\[
\mathtt{numerical\_failure}(\mathtt{empty\_topic}).
\]
此外，若$C_{k\cdot}>0$但某个活动词的$C_{kv}=0$，则命名为
\[
\mathtt{numerical\_failure}(\mathtt{support\_collapse}).
\]
四条分支都保留上一合法状态，不把$\log0$或$0/0$送入数值计算。
\end{chapterexercisesolution}
\fi

\Needspace{6\baselineskip}
% EXERCISE-SOURCE: lecture_supplement
% SOLUTION-SOURCE: lecture_supplement
\begin{exercise}\label{ex:V5-C06-S04-04}
取$\alpha=(1,1)$，三词责任度为$(0.8,0.2)$、$(0.3,0.7)$、$(0.6,0.4)$。完成一次$\boldsymbol\gamma$更新；若下一轮责任度不变，说明停止判据还应检查什么。
\end{exercise}
\ifSLFullEdition
\phantomsection\label{sol:V5-C06-S04-04}
\begin{chapterexercisesolution}{ex:V5-C06-S04-04}
\nopagebreak[4]
责任度列和为$(1.7,1.3)$，故
\[
\boldsymbol\gamma
=\boldsymbol\alpha+\sum_{n=1}^{3}\boldsymbol\eta_n
=(1,1)+(1.7,1.3)=(2.7,2.3).
\]
下一轮责任度不变意味着坐标固定，但停止还须核验
\[
\frac{|\mathcal L^{(t+1)}-\mathcal L^{(t)}|}
{1+|\mathcal L^{(t)}|}<\varepsilon,\qquad
\min_k\gamma_k>0,\qquad
\max_n\left|\sum_k\eta_{nk}-1\right|<\varepsilon_{\rm norm},
\]
并要求局部ELBO有限，以排除下溢造成的伪不变。
\end{chapterexercisesolution}
\fi

\Needspace{6\baselineskip}
% EXERCISE-SOURCE: lecture_supplement
% SOLUTION-SOURCE: lecture_supplement
\begin{exercise}\label{ex:V5-C06-S04-05}
某主题对三个词的期望计数依次为$2.5,1,0.5$。求无先验的
$\boldsymbol\varphi_k$更新，说明原计数为何不是概率；再说明Newton更新
$\boldsymbol\alpha$时如何保护正性。
\end{exercise}
\ifSLFullEdition
\phantomsection\label{sol:V5-C06-S04-05}
\begin{chapterexercisesolution}{ex:V5-C06-S04-05}
\nopagebreak[4]
总计数为4，所以
\[
\boldsymbol\varphi_k
=\frac{(2.5,1,0.5)}{2.5+1+0.5}
=(0.625,0.25,0.125),\qquad
\sum_v\varphi_{kv}=1.
\]
原计数向量和为4而非1，不能作为类别概率。Newton候选先解
\[
H(\boldsymbol\alpha)\boldsymbol\Delta
=\boldsymbol g(\boldsymbol\alpha),\qquad
\boldsymbol\alpha^{+}
=\boldsymbol\alpha-\rho\boldsymbol\Delta,\quad
\rho\in\{1,\tfrac12,\tfrac14,\ldots\},
\]
再回溯到
\[
\min_k\alpha_k^{+}>0,\qquad
\mathcal L(\boldsymbol\alpha^{+})\ge\mathcal L(\boldsymbol\alpha).
\]
若步长低于下限仍失败，则保留旧值并返回诊断。
\end{chapterexercisesolution}
\fi

\Needspace{6\baselineskip}
% EXERCISE-SOURCE: lecture_supplement
% SOLUTION-SOURCE: lecture_supplement
\begin{exercise}\label{ex:V5-C06-S05-02}
为同一语料设计Gibbs与变分EM的公平比较：固定训练/验证/测试切分、$K$、词表和预算，列出至少四项诊断与两个评价指标，并说明文档完成式困惑度如何避免信息泄漏。
\end{exercise}
\ifSLFullEdition
\phantomsection\label{sol:V5-C06-S05-02}
\begin{chapterexercisesolution}{ex:V5-C06-S05-02}
\nopagebreak[4]
先冻结同一分层切分、词表、未知词规则、$K$与总计算预算：
\[
(\mathcal D_{\rm tr},\mathcal D_{\rm va},\mathcal D_{\rm te},
\mathcal V,K,B)_{\rm Gibbs}
=(\mathcal D_{\rm tr},\mathcal D_{\rm va},\mathcal D_{\rm te},
\mathcal V,K,B)_{\rm VI}.
\]
Gibbs至少报告链间主题匹配后的计数稳定性、有效样本量/自相关、留样数和随机种子；变分至少报告ELBO轨迹、不同初始化的离散程度、局部未收敛文档数和边界/空主题。
共同指标可用文档完成式留出困惑度与主题稳定/语义连贯度，并同时报告运行时间和最大工作存储量。
文档完成时只用文档的观察部分推断
\[
\widehat{\boldsymbol\theta}_m
=\operatorname{Infer}(\boldsymbol w_m^{\rm obs};\widehat{\boldsymbol\varphi}),
\qquad
\ell_m^{\rm hold}
=\sum_{v\in\boldsymbol w_m^{\rm hold}}
\log\sum_k\widehat\theta_{mk}\widehat\varphi_{kv},
\]
再以
\[
\operatorname{PPL}
=\exp\!\left\{-\frac{\sum_m\ell_m^{\rm hold}}
{\sum_m|\boldsymbol w_m^{\rm hold}|}\right\}
\]
评价被遮住的词；若遮住词参与E步，就发生信息泄漏。
\end{chapterexercisesolution}
\fi

\end{SLChapterExercisePairSection}
% KNOWLEDGE-PLAN: KN-V5-C35-CONCEPT-010 L1
\paragraph{读前自检：本章结论与下一步。}潜在狄利克雷分配章末由“LDA以文档主题向量与主题词向量形成两层混合生成模型，词袋交换性不等于词元独立”落到“折叠Gibbs执行“减—算—抽—加”，变分EM在局部责任度与全局参数之间交替”；请把前者产出和后者输入逐项对齐，固定核对前者末项的输出类型能否作为后者首项的输入类型。\par

\begin{SLCompactClosure}
\phantomsection\label{struct:V5-C06-CH36}
\SLChapterFiveSummary{LDA以文档主题向量与主题词向量形成两层混合生成模型，词袋交换性不等于词元独立。}
{两组严格正Dirichlet先验经多元Beta积分得到折叠联合；ELBO恒等式和KL差给出可归一化坐标最优。}
{折叠Gibbs执行“减—算—抽—加”，变分EM在局部责任度与全局参数之间交替。}
{适用于计数语料与明确词表；燃烧期不是收敛证明，ELBO不等于真实证据，困惑度也不等于解释性。}
{\NextChapter{36}{\PageRank{}算法}；\GlobalChapterRef{37}{5}{8}统一比较矩阵分解、概率主题模型与采样推断。}

\phantomsection\label{struct:V5-C06-CH37}
\begin{selfcheckbox}
\begin{enumerate}[leftmargin=2.2em]
  \item \textbf{概念：}能否区分$\theta/\varphi$、$\alpha/\beta$、$\eta/\varphi$，并解释交换性不等于独立？不能则回看定义\ref{def:V5-C06-lda}和\cref{sec:V5-C06-S05}。
  \item \textbf{推导：}能否不跳步推出折叠联合、留一比值、KL--ELBO恒等式和两条平均场更新？不能则回看\cref{sec:V5-C06-S03,sec:V5-C06-S04}的正式结果。
  \item \textbf{算法：}能否写出Gibbs“减--算--抽--加”顺序，以及变分EM内外循环、停止条件和异常分支？不能则回看算法\ref{alg:V5-C06-gibbs}--\ref{alg:V5-C06-vem}。
  \item \textbf{应用与复现：}能否为新语料冻结切分、词表和预处理，选择$K$，设计多启动/多链，并在相同初始化预算下用无泄漏困惑度、稳定性和诊断公平比较？不能则回看\cref{sec:V5-C06-S05}。
  \item \textbf{诊断：}能否区分链未混合、局部ELBO未收敛、空主题与评价泄漏，并为每类问题给出对应证据？不能则回看\cref{sec:V5-C06-S05}的诊断要求；五类均能独立完成，才算掌握本章，只会调用现成主题模型而不能解释计数、ELBO和评价切分，不视为通过。
\end{enumerate}
\end{selfcheckbox}
\end{SLCompactClosure}
